%%% historicity.tex — master document
%%% Compile: pdflatex historicity && bibtex historicity && pdflatex historicity && pdflatex historicity
%%%
\documentclass[12pt,oneside]{article}

%% –– Geometry & typography ––
\usepackage[letterpaper,margin=1in,headheight=14pt]{geometry}
\usepackage[T1]{fontenc}
\usepackage[utf8]{inputenc}
% microtype omitted — requires scalable fonts (lmodern not available)
\usepackage{amsmath,amssymb}
\usepackage{setspace}
\onehalfspacing

%% –– Sectioning & headers ––
\usepackage{titlesec}
\usepackage{fancyhdr}
\pagestyle{fancy}
\fancyhf{}
\fancyhead[L]{\small\itshape The Evidentiary Status of the Historical Jesus}
\fancyhead[R]{\small\thepage}
\renewcommand{\headrulewidth}{0.4pt}

%% –– References & links ––
\usepackage[authoryear,sort&compress]{natbib}
\bibliographystyle{plainnat}
\usepackage[colorlinks=true,linkcolor=blue!60!black,
            citecolor=blue!60!black,urlcolor=blue!60!black]{hyperref}

%% –– Miscellaneous ––
\usepackage{booktabs}
\usepackage{enumitem}
\IfFileExists{csquotes.sty}%
  {\usepackage[autostyle,english=american]{csquotes}}%
  {\providecommand{\enquote}[1]{``##1''}}
\usepackage[dvipsnames]{xcolor}
\IfFileExists{epigraph.sty}%
  {\usepackage{epigraph}}%
  {\providecommand{\epigraph}[2]{%
     \begingroup\par\hfill\begin{minipage}{0.75\linewidth}%
     \itshape\small ##1\par\vspace{0.3em}\normalfont\hfill---~##2%
     \end{minipage}\par\vspace{1em}\endgroup}}

%% –– Greek text support ––
\IfFileExists{textgreek.sty}{\usepackage{textgreek}}{}   % optional; all Greek in this document is transliterated

%% –– Custom commands ––
\newcommand{\eg}{\textit{e.g.}}
\newcommand{\ie}{\textit{i.e.}}
\newcommand{\cf}{\textit{cf.}}
\newcommand{\circa}{\textit{c.}\,}
\newcommand{\grk}[1]{\textit{#1}}   % transliterated Greek

%% ============================================================
%%  DOCUMENT
%% ============================================================
\begin{document}

%% –– Title block ––
\begin{center}
  {\LARGE\bfseries The Evidentiary Status of the Historical Jesus:\\[4pt]
   A Critical Assessment}\\[18pt]
  {\large Stan Thomas}\\[4pt]
  {\normalsize Independent Researcher, St.\ George, Utah\\
   ORCID: 0000-0002-8828-7856}\\[12pt]
  {\normalsize Draft of \today}
\end{center}

\vspace{1em}

%% –– Abstract ––
\begin{abstract}
\noindent
The historicity of Jesus of Nazareth is treated as settled in mainstream
New Testament scholarship, with near-universal consensus that a specific
historical individual underlies the Christian tradition.  This paper
subjects that consensus to an evidence-first examination, evaluating
each category of source material on its own terms: the Pauline
epistles, the Synoptic and Johannine Gospels, non-canonical Christian
texts, the non-Christian references in Josephus and Tacitus, and the
documentary record.  I argue that (1)~the Pauline corpus primarily
describes a revelatory, celestial soteriological figure assembled from
pre-Christian Jewish apocalyptic and mystical sources---the Danielic
Son of Man, the Suffering Servant, Philo's Logos, angelomorphic
mediators---rather than from pagan mystery religions; a small number
of Pauline passages nonetheless admit historical readings
(Gal~4:4, Rom~1:3, 1~Cor~15:3--8) that the mythic-origin model must
and does not fully resolve.  (2)~The Gospel narratives are
demonstrably constructed from Hebrew Bible templates; this does not
prove the non-existence of a historical subject, but it does
eliminate the Gospels as independent testimony for such a subject,
shifting the evidentiary burden to non-Gospel strata.  (3)~Non-
Christian attestation is weaker than typically claimed; it does,
however, confirm the early existence of Christian traditions about a
historical founder, which is datum both hypotheses must accommodate.
(4)~The conspicuous absence of contemporaneous documentation across
multiple independent traditions constitutes disconfirming evidence
whose cumulative weight is qualitatively real but difficult to
quantify precisely.  (5)~The source trajectory from mythic to
historical is the reverse of what the standard model predicts, though
the speed of the Pauline-to-Markan transition remains a challenge
for the mythic-origin scenario that a rigorous version of it must
engage rather than dismiss.  (6)~The single passage most frequently
cited as decisive for historicity---Galatians~1:19, \enquote{James
the brother of the Lord}---is more ambiguous than the consensus
acknowledges when examined in its full Pauline context and against
the well-attested Greek religious-associative use of
\grk{adelphoi tou [patron]} as a community-title, as documented by
\citet{Harland2005}.  A corpus-linguistic baseline study of ~212
\textit{adelph-} tokens across six non-Christian contemporary-or-
earlier Greek corpora accompanies the paper as supplementary
material.  The evidence is insufficient to establish the historicity
of Jesus with the confidence the consensus claims, and the
mythic-origin model provides a parsimonious account of much of the
evidence.  The responsible conclusion is that the question is more
open than mainstream scholarship admits, that the mythicist
hypothesis---particularly in the form developed by
\citet{Carrier2014} which draws on Jewish rather than pagan
apocalyptic resources---deserves serious engagement rather than
dismissal, and that the scholarly consensus is sustained partly by
institutional and methodological inertia rather than by evidential
weight alone.
\end{abstract}

\vspace{0.5em}
\noindent\textbf{Keywords:} historical Jesus, mythicism, Pauline
epistles, Gospel sources, criterion of embarrassment, dying-and-rising
gods, Josephus, \grk{adelphos tou kyriou}, euhemerization,
historiography

\newpage
\tableofcontents
\newpage

%% ============================================================
%%  SECTIONS — to be added incrementally
%% ============================================================

\section{Introduction: The Question and Its Stakes}
\label{sec:intro}

\epigraph{Extraordinary claims require extraordinary evidence, but so do
ordinary claims when the ordinary evidence is missing.}{}

Did a specific historical individual named Jesus of Nazareth---a Jewish
preacher from Galilee, crucified under Pontius Pilate, whose followers
founded the movement that became Christianity---actually exist?  The
question is routinely treated as settled.  Virtually every survey of the
field, from \citet{Schweitzer1906} through \citet{Meier1991} to
\citet{Ehrman2012}, asserts that the historicity of Jesus is as certain
as any fact of the ancient world.  \citet{Ehrman2012} writes that
``virtually every competent scholar of antiquity, Christian or
non-Christian, agrees'' that Jesus existed.  \citet{Dunn2003} declares
that ``today nearly all historians, whether Christians or not, accept
that Jesus existed.''

This paper asks a simple question: does the evidence warrant that
confidence?

The answer developed here is that it does not.  When each category of
evidence is examined on its own terms---without deference to the
consensus as a prior, without assuming what needs to be demonstrated, and
without suppressing the evidential weight of absent evidence---the
positive case for a specific historical individual is remarkably thin.
It reduces, on close examination, to the interpretation of a single
ambiguous phrase in a single letter of Paul, supported by heuristic
arguments about what is typical of religious movements.  Everything else
that is routinely cited in support of historicity either fails textual
scrutiny, derives from the internal tradition of the movement itself, or
is indistinguishable from what we would expect if no historical
individual existed.

This is not a polemical conclusion.  It is not motivated by hostility to
religion or by allegiance to any school of thought.  It is the result of
applying to this particular question the same evidential standards that
historians apply to every other question about the ancient world---and
of refusing to grant the historicity of Jesus a special epistemic
exemption from those standards.

The stakes are significant.  If the historical Jesus is not established
by the evidence, then the entire discipline of ``historical Jesus
research''---three successive ``quests'' spanning two centuries of
scholarship \citep{Schweitzer1906,Kasemann1954,Wright1996}---has been
investigating a figure whose existence it presupposed rather than
demonstrated.  The reconstructions that have emerged from these
quests---Schweitzer's apocalyptic prophet, Crossan's peasant cynic
\citep{Crossan1991}, Sanders's eschatological preacher
\citep{Sanders1985}, Wright's self-conscious Messiah
\citep{Wright1996}, Ehrman's apocalyptic herald
\citep{Ehrman1999}---would then be understood not as competing portraits
of a historical person but as competing exercises in retrojecting
scholarly assumptions into an evidentiary void.  This does not make the
Christian tradition uninteresting or unimportant; it does, however,
fundamentally change what it is and how it should be studied.

The structure of the paper is as follows.  Section~\ref{sec:method}
establishes the methodological framework: the distinction between
evidence and authority, the Bayesian structure of historical
inference, the formal conditions under which absence of evidence
constitutes evidence of absence, the collapse of the criteria of
authenticity, and a comparative analysis of how existence-questions
have been resolved for other contested historical figures
(Pilate, Socrates, Moses, Lycurgus, Pythagoras, Romulus, Arthur,
Homer, Laozi).  Section~\ref{sec:paul} examines the Pauline corpus,
the earliest stratum of Christian writing, with particular attention
to the dating and textual-critical state of contested passages
(Gal~1:19, Rom~1:3--4, 1~Cor~15:3--8, Gal~4:4, 1~Thess~2:14--16)
and to recent Paul-skeptical scholarship.
Section~\ref{sec:gospels} analyses the Gospel sources as literary
compositions built from identifiable scriptural templates, with
explicit attention to the logical step from ``Gospels are not
independent historical testimony'' to ``a historical figure is
neither proved nor disproved.''  Section~\ref{sec:noncanon}
extends the analysis to non-canonical Christian texts.
Section~\ref{sec:external} evaluates the non-Christian attestation
channels: the Josephan \textit{Testimonium} and James passages,
Tacitus, Pliny and Suetonius, the rabbinic and \textit{Toledoth
Yeshu} traditions, the Mandaean tradition, and the archaeological
substrate (Nazareth habitation, Sepphoris silence, the Yehohanan
ossuary, the James ossuary controversy).
Section~\ref{sec:silence} develops the argument from silence across
multiple independent documentary channels.  Section~\ref{sec:myth}
presents the prior-sources argument, locating the Jesus tradition's
material primarily in the Jewish apocalyptic-mystical matrix
(Danielic Son of Man, Suffering Servant, angelomorphic mediators,
Philo's Logos), engaging the revelatory-ascent mechanism through
which Paul himself claims to have received his Christ-knowledge,
and treating Hellenistic parallels as secondary ambient context
rather than as primary source.  Section~\ref{sec:trajectory}
examines the trajectory of source development and argues that it
runs from mythic to historical, not the reverse.
Section~\ref{sec:sociology} considers the sociological structure of
the consensus and engages with the most rigorous historicist
interlocutors (Ehrman, Casey, Meier, Allison).
Section~\ref{sec:origins} offers a constructive mythic-origin
scenario, with honest acknowledgment of its unresolved weak points.
Section~\ref{sec:conclusion} summarizes the findings.


%% ============================================================
\section{Methodological Preliminaries}
\label{sec:method}
%% ============================================================

Before examining the evidence, it is necessary to establish the
principles by which evidence will be evaluated.  Much of the
disagreement in the historicity debate arises not from different
assessments of the same evidence but from different---and often
unstated---methodological commitments about what counts as evidence, how
much weight should be given to scholarly consensus, how absence of
evidence should be treated, and what the appropriate framework for
probabilistic historical reasoning looks like.

%% ------------------------------------------------------------
\subsection{Evidence, Authority, and Consensus}
\label{sec:method:authority}
%% ------------------------------------------------------------

Arguments from authority---including arguments from scholarly
consensus---are heuristics, not evidence.  They are useful when one
lacks the time, expertise, or access to primary sources needed to
evaluate a question independently.  But they are defeasible: when the
primary evidence is examined and found to be in tension with the
consensus, the evidence takes precedence.  This is not a controversial
epistemological principle; it is the foundation of all empirical inquiry.

In the case of the historical Jesus, the consensus is routinely invoked
not as a heuristic but as evidence in its own right.  Ehrman's frequent
assertion that ``virtually all scholars'' agree is presented as though
the agreement itself constitutes a reason to believe, independent of the
evidence on which it rests \citep{Ehrman2012}.  But the epistemic value
of a consensus depends entirely on the independence and quality of the
evidential judgments that underlie it.  A consensus among scholars who
have all been trained within the same paradigm, who have all absorbed
the same default assumptions during their formation, and who operate
within institutions that presuppose the conclusion in question, carries
far less evidential weight than a consensus formed by independent
evaluation of primary data.

The comparison with the scientific consensus on climate change, which
Ehrman implicitly invites, is instructive precisely because it fails.
The climate consensus rests on thousands of independent empirical
measurements---ice cores, satellite data, ocean temperatures, atmospheric
chemistry---converging on the same conclusion from radically different
evidential pathways.  The historicity consensus rests on a small number
of ancient texts, all of which are internal to the tradition in question
or dependent upon it, interpreted within a framework that has
presupposed historicity since its inception.  The two situations are not
analogous.

This paper therefore treats the consensus as a datum to be explained,
not as evidence to be weighed.  The question is not whether scholars
agree but whether the evidence warrants their agreement.

%% ------------------------------------------------------------
\subsection{The Bayesian Framework}
\label{sec:method:bayes}
%% ------------------------------------------------------------

Historical inference is inherently probabilistic.  We cannot observe the
past directly; we can only assess the probability of competing
hypotheses given the surviving evidence.  The appropriate formal
framework for this assessment is Bayesian: we compare the probability of
the evidence under each hypothesis and update our credence accordingly
\citep{Carrier2012}.

Let $H$ denote the hypothesis that a specific historical individual
underlies the Christian tradition, and let $M$ denote the hypothesis
that the Jesus figure originated as a mythic-revelatory construct
subsequently historicized.  For any piece of evidence $E$, Bayes'
theorem gives:

\begin{equation}
\label{eq:bayes}
\frac{P(H \mid E)}{P(M \mid E)}
  = \frac{P(E \mid H)}{P(E \mid M)} \cdot \frac{P(H)}{P(M)}
\end{equation}

\noindent
The ratio $P(E \mid H) / P(E \mid M)$ is the \textit{likelihood ratio}:
how much more (or less) expected the evidence is under historicity than
under mythicism.  The ratio $P(H)/P(M)$ represents our prior
assessment before considering that evidence.

\citet{Carrier2014} applied this framework systematically in \textit{On
the Historicity of Jesus}, arriving at a posterior probability for
historicity in the range of roughly one in three.  The specific numbers
Carrier assigns to individual priors and likelihoods are contestable,
and \citet{Vridar2014} have raised legitimate concerns about circularity
in the construction of the mythicist model's ``predictions.''  But the
framework itself is sound, and it has the virtue of making explicit what
is usually left implicit: every historical judgment involves an
assessment of priors, an assessment of likelihoods, and an update.
Scholars who reject the Bayesian framing are not avoiding these
assessments; they are making them informally and without accountability.

Two features of the Bayesian framework are particularly relevant here.
First, evidence that is equally probable under both hypotheses has no
discriminating power, regardless of how impressive it may seem in
isolation.  If a piece of evidence is exactly what we would expect
whether or not Jesus existed---\eg, the existence of early Christian
communities, or the presence of a plausible Palestinian setting in the
Gospels---then it cannot be cited in favor of either hypothesis.  Much
of what is routinely presented as evidence for historicity fails this
test.

Second, the cumulative force of evidence matters.  Individual pieces of
evidence may each have modest likelihood ratios, but if they
consistently favor one hypothesis over the other, the cumulative effect
can be decisive.  This is particularly important for the argument from
silence developed in Section~\ref{sec:silence}: no single absent source
is conclusive, but the systematic absence of expected evidence across
multiple independent traditions produces a cumulative likelihood ratio
that substantially favors the mythic-origin hypothesis.

Third, \citet{Tucker2004} has noted that Bayesian historical inference
is particularly sensitive to the independence of the testimonies on
which likelihood estimates rest: the cumulative case for any historical
claim requires that the contributing sources be independently probative,
and textual-critical research showing that Gospel attestations collapse
to Markan or Q dependence (Section~\ref{sec:gospels:q}) accordingly
reduces the effective likelihood contribution rather than multiplying
it.  This is an unglamorous technical point, but it bears directly on
the strongest traditional historicist argument---multiple independent
attestation---and its implications have been under-pursued in the
historical-Jesus literature.

\paragraph{A sensitivity analysis implementation.}
The project repository includes a simple dependency-aware Bayesian
sensitivity analysis (\textit{scripts/bayesian\_sensitivity.py},
report in \textit{baseline/bayesian\_report.md}) that implements
this framework explicitly for the present question.  Each of the
principal evidentiary strands discussed in this paper
(Paul's \textit{adelphos} distribution, Rom~1:3, Gal~4:4, the
1~Corinthians~15 witness list, Paul's biographical silence, the
Gospels' scriptural construction, the Josephan James passage, the
Testimonium Flavianum, Tacitus, the Jewish-apocalyptic matrix,
cumulative silence) is assigned a point-estimate likelihood ratio
in decibels, with source-dependency explicitly flagged so that the
contribution of dependent sources is discounted rather than
multiplied.  Under plausible point estimates and a neutral
$P(H)=0.5$ prior, the posterior $P(H|E)$ lands in the range of
roughly 0.23--0.39 depending on how aggressively dependencies are
discounted---open-question territory by any reasonable standard.
Under higher priors ($P(H)=0.9$), the posterior remains higher
(0.73--0.85) but still materially below the near-certainty the
historicity consensus projects.

The tool is explicitly a sensitivity analysis rather than a
claim to the ``correct'' posterior: point estimates are
contestable, and the exercise is to see how the posterior shifts
under different assumptions about individual evidence pieces.  The
tool's principal contribution is to make the informal Bayesian
reasoning of the field auditable: Carrier's \citeyear{Carrier2014}
multiplication approach, extended or contested point-by-point by
historicist critics, can be made to run on a shared apparatus
where the sources of posterior-sensitivity are explicit.  This is
what Tucker 2004 advocates methodologically; to our knowledge it
has not been done for the historicity question until now.

%% ------------------------------------------------------------
\subsection{Absence of Evidence as Evidence of Absence}
\label{sec:method:silence}
%% ------------------------------------------------------------

The aphorism ``absence of evidence is not evidence of absence'' is
widely repeated and, in many contexts, flatly wrong.  The correct
principle, recognized in Bayesian epistemology and in standard
scientific reasoning, is that absence of evidence \textit{is} evidence
of absence precisely when the evidence would be expected to exist if the
hypothesis were true \citep{Sober2009,Carrier2012}.

Formally: if $P(E \mid H) > P(E \mid \neg H)$, then observing $E$
supports $H$; and by the same logic, if we would expect to observe
evidence $E$ given hypothesis $H$ but we do not observe it, then the
absence of $E$ is evidence against $H$.  The strength of this evidence
is proportional to how confidently we would expect $E$ under $H$.

Consider an analogy.  If someone claims that a major fire destroyed a
city block last Tuesday, and we find no fire department records, no
insurance claims, no news reports, no physical damage, and no
eyewitness accounts---despite all of these being the kinds of evidence
that reliably exist when major fires occur---the absence of all this
expected evidence is strong evidence that the fire did not happen.  We
do not say ``absence of evidence is not evidence of absence'' and
conclude that the question is open.  We conclude, reasonably, that the
claim is almost certainly false.

The application to the historical Jesus is direct.  First-century Roman
Palestine is one of the better-documented periods and regions of the
ancient world.  We possess extensive writings by Philo of Alexandria,
covering Jewish theology, politics, and interactions with Roman
authority in exactly the relevant period.  We know of a history of
Galilee by Justus of Tiberias, covering Jesus's supposed home region.
The Mishnaic traditions preserve detailed discussions of Sanhedrin
procedure, heresy, and religious authority.  Roman administrative
practices generated records of trials and executions.  If a figure
performed even a modest fraction of what the Gospels describe---public
ministry attracting large crowds, dramatic confrontation with temple
authorities, trial before the Sanhedrin and the Roman prefect, execution
accompanied by extraordinary portents---we would expect traces in these
independent documentary traditions.

The traces are entirely absent.  This is not merely ``an argument from
silence.''  It is an argument from the specific, predicted,
well-characterized absence of evidence that should exist if the
historicist hypothesis is true, across multiple independent channels.
In any other historical investigation, this pattern would be treated as
significant disconfirming evidence.  That it is minimized in Jesus
studies---typically with the assertion that Jesus was ``too obscure'' to
attract notice---is itself a methodological anomaly requiring
explanation.  The ``obscurity'' defense, moreover, contradicts the
Gospel narratives themselves, which describe a figure who attracted
enormous public attention.  One cannot simultaneously maintain that
Jesus was prominent enough to warrant the Gospel descriptions and
obscure enough to explain the silence of every non-Christian source.

The formal point bears emphasis: we are not asking whether a specific
piece of evidence \textit{proves} nonexistence.  No single absent
source does.  We are asking what the \textit{cumulative} likelihood
ratio is when every expected evidentiary channel returns null.  If there
are $n$ independent channels, each of which would produce evidence with
probability $p$ given historicity but probability $q < p$ given
mythic origins, the cumulative likelihood ratio is $(q/p)^n$, which
can become very small even when each individual ratio $q/p$ is only
modestly below unity.  This cumulative structure is developed in detail
in Section~\ref{sec:silence}.

%% ------------------------------------------------------------
\subsection{The Criteria of Authenticity}
\label{sec:method:criteria}
%% ------------------------------------------------------------

Historical Jesus scholarship has traditionally relied on a set of
``criteria of authenticity'' to identify material in the Gospels that
plausibly derives from a historical Jesus.  The most prominent are:

\begin{itemize}[leftmargin=2em]
\item \textbf{The criterion of embarrassment:} material that would be
embarrassing to the early church is likely historical, since the church
would not have invented it.  (Applied to the crucifixion, the baptism
by John, Peter's denial, etc.)

\item \textbf{The criterion of multiple attestation:} material attested
in multiple independent sources is more likely historical.

\item \textbf{The criterion of dissimilarity:} material that is
unlike both Judaism and later Christianity is likely to derive from the
historical Jesus himself, since it cannot be attributed to either the
Jewish context or the later church.

\item \textbf{The criterion of coherence:} material that coheres with
already-established authentic material is also likely authentic.
\end{itemize}

These criteria have been subjected to devastating criticism from within
mainstream scholarship itself.  The criterion of embarrassment assumes
we can reliably determine what the early church would have found
embarrassing---an assumption that requires knowledge of early Christian
psychology we do not possess (Section~\ref{sec:myth:crucified}).  The
criterion of multiple attestation collapses if the supposedly
independent sources are not in fact independent---as the Synoptic
problem demonstrates (Section~\ref{sec:gospels:q}).  The criterion of
dissimilarity is methodologically absurd: it can only identify a Jesus
who was neither Jewish nor Christian, which is historically incoherent.
The criterion of coherence is circular: it identifies material that
fits a portrait already constructed by the other criteria.

Recent scholarship has increasingly acknowledged these problems.
\citet{Keith2024} and collaborators have explicitly argued that ``the
era of the hard version of the criteria of authenticity is over,'' and
that the field must move toward memory theory and reception history.
\citet{Allison2010} has expressed deep skepticism about the criteria's
ability to recover specific historical sayings or events.  Even
\citet{Ehrman2012}, while defending historicity, relies less on the
criteria than on broader arguments about source plurality and the
general character of the tradition.

The collapse of the criteria of authenticity within mainstream
scholarship is significant for the present inquiry: it means that the
principal tools by which historicists have claimed to identify a
``historical core'' beneath the theological overlay are now recognized
as methodologically unreliable by the historicists themselves.  If the
criteria cannot reliably distinguish historical material from
theological invention, the Gospels cannot serve as evidence for a
historical figure---they can only serve as evidence for the beliefs and
literary practices of the communities that produced them.

%% ------------------------------------------------------------
\subsection{The Existence Question in Comparative Perspective}
\label{sec:method:comparative}
%% ------------------------------------------------------------

Discussions of the historicity of Jesus are routinely conducted as if
the question were unique.  It is not.  A class of comparable figures
exists whose putative historicity has been seriously questioned within
modern scholarship and in many cases definitively resolved.  Examining
how these cases have been handled yields diagnostic criteria that can
be applied to Jesus's evidentiary profile.  This comparative exercise
is not argument by analogy---each case has its own particulars---but a
methodological check: we can ask what kinds of evidence have
historically been \textit{decisive} in resolving existence-questions,
and how Jesus's evidentiary profile stands against those criteria.

\paragraph{The class of figures.}
At least three categories are relevant.

\textit{Category I: Confirmed historical, formerly doubted.}
\textbf{Pontius Pilate} was, until 1961, known only from literary
sources (Philo, Josephus, Tacitus, Gospels) and his historicity had
occasionally been questioned on those grounds.  The discovery of the
Caesarea Maritima inscription---a dedication stone naming Pilate as
``prefect of Judaea''---settled the question definitively
\citep{Vardaman1962}.  Here, physical epigraphic evidence
contemporary with the figure's lifetime was decisive.
\textbf{Socrates} was a figure known only through literary sources
until recent times (and still chiefly so), but multiple
\textit{independent} contemporary witnesses---Plato (sympathetic),
Xenophon (admiring), Aristophanes (hostile, \textit{Clouds} 423~BCE),
and Aristotle (who cites teachers of teachers)---converge on a
specific individual with specific biographical details.  Convergent
independent attestation from living contemporaries, including hostile
and satirical ones, is a distinguishing feature.
\textbf{Paul} himself is historically secure: we possess letters
written by him and addressed to specific communities, quoted by
second-century Christian authors, and attested in the archaeological
and epigraphic record of the specific cities he addresses.

\textit{Category II: Rejected as historical.}
\textbf{Moses} is treated as non-historical by a significant majority
of modern Egyptologists and Hebrew Bible scholars.  The decisive move
was negative: the \textit{absence} of archaeological or documentary
evidence in Egypt for a mass Hebrew migration, coupled with the
\textit{internal} literary-critical analysis showing the Mosaic
narratives as composite and late
\citep{FinkelsteinSilberman2001,ThompsonTL1999,Dever2003}.
Here, archaeological silence in a period and place where such
evidence would be expected was determinative.
\textbf{Lycurgus of Sparta}, the legendary founder of the Spartan
constitution, is now generally treated as a non-historical figure
\citep{Cartledge2001}: the ancient sources that purport to describe
him date from centuries after his supposed lifetime, the
biographical content is schematic and founder-myth-conforming, and
his constitutional reforms are more plausibly attributed to
incremental institutional development than to an individual reformer.
Here, \textit{late and derivative literary sources coupled with
founder-myth conformity} was decisive.
\textbf{William Tell} was accepted as historical for centuries but
was shown by nineteenth-century philological scholarship to be a
folkloric reshaping of the apple-shot motif attested in Saxo
Grammaticus and in Norse tradition centuries earlier---the latest
plausible moment for a historical kernel is around 1200 CE but the
narrative has been demonstrated to be literary transmission rather
than historical memory.

\textit{Category III: Genuinely contested.}
\textbf{Pythagoras} sits between the first two categories: a
probable historical kernel (5th-century BCE, pre-Socratic,
Samos/Croton) surrounded by dense legendary accretion beginning
within decades of his death \citep{Riedweg2005}.  No contemporary
source attests him directly; our earliest references are a
generation later; most biographical detail is Platonic or
Neoplatonic reshaping.  Modern scholarship accepts a minimal
historical Pythagoras while treating most of the biographical
tradition as non-historical.
\textbf{Romulus}, the supposed founder of Rome, is now universally
rejected as historical \citep{Beard2015,Wiseman1995} despite a
sophisticated literary tradition that placed him in a specific
historical setting (753~BCE).  The \textit{absence of any
independent attestation} and the \textit{demonstrable origin of the
narrative in later Roman foundation-myth requirements} were
decisive.  Romulus is the canonical example of a figure whose
biography was \textit{generated} to satisfy a community's need for a
founder, rather than being remembered.
\textbf{King Arthur} remains contested.  The literary accretions are
extensive (Geoffrey of Monmouth, the Vulgate Cycle); the earliest
references are late (Nennius, \textit{c.}~830~CE); archaeological
evidence for a specific Arthur is absent; and most recent specialist
scholarship regards him as legendary with at most a minimal
historical kernel in post-Roman Britain \citep{Higham2002}.
\textbf{Homer} is the longest-running case: the authorship question
(``Homeric Question'') has oscillated between single-author,
multiple-author, and oral-tradition views for two centuries
\citep{WestML2011,Finkelberg2005}.  No contemporary evidence; late
and uncertain biographical tradition; the text itself is the primary
evidence and bears internal signs of composite formation.
\textbf{Laozi}, traditional founder of Daoism, is now regarded by
most specialists as either a legendary figure composed of several
traditional strands or a retrospective construction
\citep{GrahamA1986}.  The \textit{Daodejing} is attributed to him
but the text's authorship is demonstrably composite.

\paragraph{Diagnostic criteria, extracted.}
From these precedents, the following criteria have historically
carried decisive weight in resolving existence-questions:

\begin{enumerate}[leftmargin=2em]
\item \textbf{Direct contemporary physical/epigraphic attestation.}
  A contemporary inscription or physical artifact naming the figure
  settled the Pilate question.  Jesus has no direct contemporary
  physical or epigraphic attestation.
\item \textbf{Multiple \textit{independent} contemporary literary
  attestation, including hostile or satirical sources.}  Socrates
  has Plato + Xenophon + Aristophanes + indirect Aristotle.
  Jesus has Paul (one step removed, writing roughly twenty years
  after the supposed crucifixion, to communities across the
  Mediterranean, and never having met a living Jesus) and later
  Gospel material that demonstrably derives from Mark.  Non-Christian
  attestation is post-contemporary and partially or wholly dependent
  on Christian informants.
\item \textbf{Archaeological silence in a period and region where
  evidence would be expected.}  Moses is rejected partly on this
  ground.  Jesus shares the feature: the Pilate prefecture is
  documented; the crucifixion of one Galilean preacher would normally
  leave no administrative trace simply because such events were
  common, but the absence is genuine.  Section~\ref{sec:silence}
  develops the cumulative silence argument.
\item \textbf{Literary-critical evidence of composite or derivative
  origin in the putatively-biographical sources.}  Lycurgus,
  Romulus, and Laozi share this feature.  Section~\ref{sec:gospels}
  argues the Gospels share it.
\item \textbf{Demonstrable origin in community need for a founder or
  theological anchor.}  Romulus is the canonical case.  The
  Pauline-to-Markan trajectory analysed in
  Section~\ref{sec:trajectory} proposes this pattern for Jesus.
\item \textbf{Dense legendary accretion within a short time.}  This
  is ambiguous---both Pythagoras (historical kernel) and Romulus
  (no historical kernel) exhibit it.  The criterion is not
  decisive on its own.
\end{enumerate}

\paragraph{Jesus's evidentiary profile in this landscape.}
Applying these criteria:

\begin{itemize}[leftmargin=2em]
\item \textbf{Against the confirmed-historical category (Pilate,
  Socrates, Paul).}  Jesus lacks the direct epigraphic attestation
  of Pilate.  He lacks the multiple independent contemporary
  literary attestation of Socrates.  He lacks the first-person
  documentary record of Paul.  His evidentiary profile is
  significantly weaker than any of these cases.
\item \textbf{Against the rejected-as-historical category (Moses,
  Lycurgus, William Tell).}  Jesus's case is stronger than
  these.  The absence-of-evidence is not as categorical
  (Sections~\ref{sec:silence:philo},~\ref{sec:silence:mishnah});
  the community tradition is immediate rather than centuries
  later; and the Pauline stratum, while non-biographical, is
  demonstrably early.
\item \textbf{Against the contested category (Pythagoras, Romulus,
  Arthur, Homer, Laozi).}  Jesus's profile most closely matches
  this category.  Dense legendary accretion within a short time;
  literary-critical evidence of derivative origin in the
  biographical sources; absence of direct epigraphic attestation;
  partial and contested non-tradition attestation; ambiguous
  archaeological silence.  He shares Pythagoras's profile of
  ``probable historical kernel with massive legendary
  elaboration''; he shares Romulus's profile of ``community
  construction to satisfy founder-myth needs''; he shares
  Arthur's profile of ``contested minimum with no definitive
  resolution.''
\end{itemize}

\paragraph{What this implies.}
The observation is not that Jesus definitely falls into Category~II
or Category~III.  It is that the methodological framework of
\textit{modern scholarship on comparable figures} does not license
the level of confidence the historicity consensus projects.  A
figure with Jesus's evidentiary profile, if he were not the founder
of a globally significant religious tradition, would be treated as
genuinely contested rather than as obviously historical.  This paper
argues that the evidentiary standards applied to comparable cases
should be applied here also---and that when they are, the
responsible conclusion is that the question is genuinely open.

A further observation: among the Category~III figures, Jesus's case
is \textit{distinctive} in exactly one way: the textual and
communal trace his tradition has left is orders of magnitude
larger than any other contested figure's.  This is a datum, but it
is not evidence of historicity.  The sheer volume of tradition a
figure generates is a function of the religious movement that
forms around him (or his invocation), not of his historicity.  The
Buddha, Mohammed, Moses, Homer, Pythagoras, and Jesus all anchor
traditions of enormous textual and communal extent; their
historicity profiles differ dramatically.  Scale of reception is
not evidence of the scale of the underlying historical event.

%% ============================================================
\section{The Pauline Corpus}
\label{sec:paul}
%% ============================================================

The letters attributed to Paul constitute the earliest stratum of
surviving Christian writing.  If a historical Jesus existed and was known
to Paul's contemporaries, this is where we would expect to find the
clearest traces.  The Pauline corpus is therefore the most important
body of evidence in the historicity debate---and, as we shall see, the
most damaging to the historicist case.

%% ------------------------------------------------------------
\subsection{Dating, Authenticity, and Interpolation}
\label{sec:paul:dating}
%% ------------------------------------------------------------

Modern scholarship recognizes seven letters as authentically Pauline:
Romans, 1~Corinthians, 2~Corinthians, Galatians, Philippians,
1~Thessalonians, and Philemon \citep{Ehrman2004}.  The remaining six
letters in the traditional Pauline corpus (Ephesians, Colossians,
2~Thessalonians, 1--2~Timothy, Titus) are widely regarded as
pseudepigraphical---written by later authors in Paul's name.  The
authentic letters are typically dated to the 50s~CE, making them the
earliest surviving Christian documents by a margin of at least fifteen
years over the earliest Gospel (Mark, \circa 70~CE).

The ``seven authentic letters'' consensus, however, deserves more
scrutiny than it typically receives.  It rests primarily on stylometric
analysis and theological consistency arguments---methods that can
identify gross differences between letter groups but that are less
reliable at detecting interpolations \textit{within} otherwise authentic
letters.  And interpolation within the authentic letters is not
hypothetical; it is widely acknowledged.  Several passages in the
undisputed letters are recognized as probable later insertions:

\begin{itemize}[leftmargin=2em]
\item \textbf{1~Corinthians 14:34--35}, commanding women to be silent
in church, is recognized by a broad range of scholars as a
post-Pauline interpolation on textual, stylistic, and theological
grounds \citep{Payne2009,Walker2001}.  Some manuscripts place the
passage after verse~40, indicating scribal uncertainty about its
location---a hallmark of marginal glosses that migrated into the text.

\item \textbf{1~Thessalonians 2:14--16}, condemning ``the Jews'' for
killing Jesus and the prophets, contains language and theology
inconsistent with Paul's views expressed elsewhere (particularly
Romans~9--11) and appears to reference the destruction of the Temple
in 70~CE (``wrath has come upon them at last''), placing it after
Paul's death \citep{Pearson1971,Schmidt1983}.

\item \textbf{1~Corinthians 15:3--11}, the resurrection creed, is
universally recognized as a pre-Pauline formula that Paul is quoting
rather than composing \citep{Conzelmann1975}.  Its boundaries,
however, are disputed: how much of the surrounding material is Paul's
own commentary and how much is part of the received formula?  The
passage contains vocabulary and constructions not characteristic of
Paul, and the question of whether the formula as we have it has itself
been expanded by later editors cannot be definitively resolved.

\item \textbf{Romans 1:3--4}, describing Jesus as ``descended from
David according to the flesh,'' is likewise identified as a
pre-Pauline creedal fragment embedded in the letter
\citep{Jewett2007}.
\end{itemize}

The Dutch Radical school of the nineteenth century---notably
\citet{VanManen1896} and their predecessors---challenged the
authenticity of the entire Pauline corpus, arguing that the letters as
we have them are products of second-century redaction.  This position
was marginalized by the emerging consensus but never systematically
refuted on textual grounds; it was displaced rather than answered.
The most sustained modern revival is \citet{Detering1995}, whose
\textit{Der gef\"alschte Paulus} reactivated the Dutch Radical
tradition in a contemporary textual-critical register.

\paragraph{Current Paul-skeptical scholarship.}
The past fifteen years have seen a genuinely new wave of academic
work on the dating and authenticity of the Pauline corpus, along two
related tracks.  The first, associated with \citet{BeDuhn2013},
\citet{Klinghardt2015}, and especially \citet{Vinzent2014,Vinzent2023},
argues for \textit{Marcion-priority}: that Marcion's mid-second-century
versions of the Pauline letters (and of Luke) are textually closer to
the originals than the canonical versions, and that the proto-orthodox
canonical forms are expansions of Marcion's text rather than Marcion's
text being an abridgement of the canonical.  The second, advanced in
\citet{Livesey2024}, argues that the Pauline letters are deliberate
rhetorical productions of the late first to early second century
written within Roman literary convention, not situational letters
from a specific historical apostle to specific churches.  Both
positions are contested; both have earned publication at Cambridge
University Press, T\&T Clark, Peeters, and Polebridge, and must be
engaged rather than dismissed.  \citet{White2014} offers a more
moderate complement: a thorough study of how different early Christian
communities actively constructed competing \textit{images} of Paul,
which implies at minimum that the figure of Paul was heavily
contested and re-shaped across the decades during which our texts
were stabilising.  \citet{Trobisch2000} and \citet{Pervo2010} provide
further relevant apparatus on the canon-formation process.

The relevance of this body of work to the historicity question is
cautious but real.  If Marcion-priority is correct, then some
apparent historicizing phrases in the canonical Paul may be
proto-orthodox insertions made in response to Marcion, and any such
phrase that is absent from Marcion's reconstructed text becomes
interpolation-suspicious.  Gal~1:19 is, significantly, \textit{present}
in Marcion's reconstructed text and so its authenticity is not
undercut by Marcion-priority; but 1~Thess~2:14--16 (already widely
recognised as interpolated) and portions of the resurrection creed
in 1~Cor~15:3--11 are more vulnerable to redaction arguments under
this framework.  If Livesey's stronger thesis is correct, the
evidentiary status of the Pauline corpus for historicity shifts more
drastically: the letters would then testify to the community
rhetoric of the early second century rather than to the biographical
memory of a specific first-century figure, which is already closer
to the mythic-origin scenario than to the standard historicist one.

This paper does not adopt either Paul-skeptical position as a load-
bearing premise.  It works within the conservative seven-letter
authenticity consensus as a stronger starting point, and the
arguments of the following sections are constructed to be valid under
that conservative assumption.  What we register here is that the
conservative assumption is under active, credentialed academic
challenge, and that the directionality of the challenge---toward later,
more constructed, and more editorially layered texts---if anything
strengthens the line of argument this paper develops.  A rigorous
engagement with the historicity question in the present decade cannot
ignore \citet{Vinzent2023}, \citet{Klinghardt2015}, \citet{BeDuhn2013},
and \citet{Livesey2024}, and we note them explicitly here so that
readers who wish to pursue the textual basis of the Pauline evidence
beyond what is defended in this paper have their citations.

\paragraph{The minimum claim required.}
We need not adopt the full Dutch Radical, Marcion-priority, or
late-Paul position to recognize that the Pauline letters are not
pristine first-century documents.  They have passed through decades
of copying, editing, and collection, during which theologically
motivated interpolation was routine---as the pseudepigraphical letters
themselves demonstrate.  Any passage in the Pauline letters that is
cited as evidence for a historical Jesus must survive not only the
question of what it means but the question of whether it is
authentically Pauline at all, or whether it represents a later
editorial layer reflecting the historicized Jesus of the emerging
Gospel tradition.

\paragraph{Textual and patristic attestation of the contested
passages.}
The honest response to interpolation arguments is to examine the
actual manuscript and patristic record for each contested passage.
The following survey, based on the critical apparatus of
\citet{NA28}, the textual scholarship of \citet{MetzgerEhrman2005}
and \citet{Zuntz1953}, the Marcion reconstructions of
\citet{BeDuhn2013} and \citet{Klinghardt2015}, and the patristic
citations catalogued in the \textit{Biblia Patristica} series,
gives a summary state for each of the passages the paper concedes
are difficult for mythicism.

\textbf{Galatians~1:19} (\grk{ton adelphon tou kyriou}).  The clause
is attested in \textbf{P46} (the Chester Beatty Pauline papyrus,
c.~200~CE and our earliest substantial witness to Galatians), in
Sinaiticus, Vaticanus, Alexandrinus, and all major later codices.
It is present in Marcion's reconstructed text as attested by
Tertullian's \textit{Adversus Marcionem} (c.~207~CE).  Its patristic
reception is early and broad: Tertullian cites it; Jerome (c.~390)
devotes an extensive commentary to the question of what kind of
``brother'' is meant \citep{Lightfoot1869}.  \textit{The clause is
textually secure.}  This paper's interpretive argument in
§\ref{sec:paul:brother} must therefore proceed on the semantic and
contextual grounds it does, not on an interpolation hypothesis.
Textual argument is unavailable to rescue mythicism from this
passage.  We accept that.

\textbf{Romans~1:3--4} (\enquote{descended from David according to
the flesh\ldots declared Son of God\ldots by resurrection}).
Attested in \textbf{P46}, Sinaiticus, Vaticanus, and the whole
manuscript tradition without significant variant affecting the
Davidic-descent clause.  Marcion's reconstructed text, however, is
\textit{disputed} at this passage: Tertullian records that Marcion's
version lacked parts of the opening salutation but the precise
state of Romans~1:3--4 in Marcion is not fully recoverable from the
patristic record, and this has been a subject of ongoing scholarly
debate (see \citealt{BeDuhn2013} for a cautious reconstruction).
More importantly for this paper's argument, Romans~1:3--4 is
\textit{widely recognized as a pre-Pauline creedal fragment}
quoted by Paul rather than composed by him \citep{Jewett2007}; the
authenticity of its incorporation into Paul's letter is secure
(P46 attests it), but the formula itself originated elsewhere in
the earliest community and represents the theological claim-making
of that community.  Textual criticism does not rescue mythicism
from the passage either, and this paper's response in
§\ref{sec:myth:davidic} accordingly proceeds on other grounds.

\textbf{Galatians~4:4} (\enquote{born of a woman, born under the
law}).  Attested in \textbf{P46} and the whole manuscript
tradition.  Textually secure.  Discussed in §\ref{sec:myth:davidic}
and in the concluding assessment.

\textbf{1~Corinthians~15:3--11} (the resurrection creed and the
witness list).  The formulaic core (verses~3--5) is universally
recognized as pre-Pauline; the boundaries of the formula are
disputed.  \textbf{P46} contains verses~3--11 without major textual
disturbance.  However, the witness-list extension in verses~6--8
(\enquote{appeared to over five hundred brothers\ldots then to James,
then to all the apostles, last of all\ldots he appeared also to me})
has long been flagged as possibly showing signs of later editorial
expansion: the disjunction between the tight creedal form
(3--5) and the looser narrative list (6--8) has invited proposals
that the list was added by Paul or by a later hand.
\citet{Conzelmann1975} treats the expansion as Pauline but not part
of the pre-Pauline formula.  A full interpolation hypothesis for
6--8 is possible but \textit{not textually supported by variant
evidence}---which means mythicist readings must accommodate the
witness list as substantially early, not dismiss it.

\textbf{1~Thessalonians~2:14--16} (\enquote{the Jews who killed the
Lord Jesus\ldots wrath has come upon them at last}).  Unlike the
passages above, this is the one textual case where interpolation is
widely supported: the anti-Jewish language contradicts Paul's
position in Romans~9--11, the reference to \enquote{wrath has come
upon them at last} appears to presuppose the destruction of the
Second Temple in 70~CE (after Paul's death), and \citet{Pearson1971}
and \citet{Schmidt1983} developed the interpolation case in detail.
\textit{But the interpolation is not attested in any surviving
manuscript variant.}  It is an inference from internal coherence
and historical plausibility.

\paragraph{What the textual record shows.}
The honest summary is that textual-critical argument alone does
\textit{not} support interpolation hypotheses for the passages that
pose the greatest difficulty for mythicism.  Gal~1:19 is attested in
P46 and across the manuscript tradition; Rom~1:3--4 is attested
(though its Marcion-priority state is partially contested); Gal~4:4
is attested; 1~Cor~15:3--8 is attested.  The interpolations most
confidently identified in modern scholarship (1~Thess~2:14--16,
1~Cor~14:34--35) are identified on internal-coherence grounds, not
on manuscript-variant grounds.

This has a direct implication for the argument of this paper: the
mythicist response to the conceded-difficulty Pauline passages
cannot rely on textual-critical excision.  It must rely on semantic,
contextual, and historical-compositional analysis, which is where
§\ref{sec:paul:brother}, §\ref{sec:myth:davidic}, and
§\ref{sec:myth:crucified} accordingly concentrate their argumentative
weight.

What the textual record \textit{does} show is that the Pauline
corpus was stable by the late second century (the P46 era), that
its pre-second-century state is accessible only through Marcion
reconstructions and patristic citation, and that some passages
widely accepted as interpolations cannot be detected from the
manuscript record alone---which is a methodological caution for
both sides of the debate.

\paragraph{The patristic citation record for the contested
passages.}
An independent check on textual stability is to ask which of the
contested passages are cited or paraphrased by the earliest
extant Christian writers outside the New Testament---the
Apostolic Fathers, writing c.~95--150~CE, some of them within the
same generation as the latest Pauline letters.  A systematic
examination of the Lightfoot collection of the Apostolic Fathers
yields a striking asymmetry.

\textbf{Romans 1:3--4} (the Davidic-descent creed) is cited by
Ignatius of Antioch (writing c.~107--110~CE) repeatedly and in
multiple letters: \textit{Ephesians} 18.2 (\enquote{Jesus the
Christ\ldots of the seed of David}), \textit{Ephesians} 20.2
(\enquote{Jesus Christ, who after the flesh was of David's
race}), \textit{Trallians} 9.1 (\enquote{Jesus Christ, who was
of the race of David, who was the Son of Mary, who was truly
born}), \textit{Romans} 7.3 (\enquote{flesh of Christ who was of
the seed of David}), and \textit{Smyrnaeans} 1.1 (\enquote{our
Lord\ldots is truly of the race of David according to the
flesh}).  In Ignatius's anti-docetic polemic, these citations are
evidentially load-bearing: they are the specific Pauline material
he adduces to establish Jesus's true embodiment against docetic
opponents.  Rom 1:3--4 is therefore demonstrably stable in the
Pauline textual tradition by c.~110~CE.

\textbf{Galatians 4:4} (\enquote{born of a woman}) is echoed in
Ignatius, \textit{Smyrnaeans} 1.1 (\enquote{truly born of a
virgin}), again in an anti-docetic context.

\textbf{Galatians 1:19} (\enquote{brother of the Lord}) is, by
contrast, \textit{not cited or paraphrased by any of the
Apostolic Fathers}.  The name \enquote{James} does not appear in
the collection.  The phrase \enquote{brother of the Lord} appears
nowhere.  First-century and early second-century Christian
writers whose works survive do not cite the passage that modern
historicist arguments treat as virtually decisive for
establishing the historicity of Jesus.  The first clear patristic
citations of the phrase come with Tertullian's \textit{Adversus
Marcionem} (c.~207~CE) and Jerome's \textit{Adversus Helvidium}
(c.~390~CE), whose extended discussion distinguishes three
patristic positions on what kind of \enquote{brother} is meant
\citep{Lightfoot1869}.

\textbf{1 Corinthians 15:3--8} (the resurrection creed and witness
list) is also not cited in the Apostolic Fathers.  1~Clement~47.1
explicitly knows 1~Corinthians as a letter and cites its
opening chapter (\enquote{concerning himself and Cephas and
Apollos}, paralleling 1~Cor~1:12), but does not cite the
resurrection witness list.  Polycarp and Ignatius likewise know
Paul's corpus but do not cite 1~Cor~15:3--8.

\paragraph{What the patristic silence does and does not imply.}
The absence of Apostolic Fathers citations for Gal~1:19 and 1~Cor
15:3--8 does \textit{not} establish that these passages were
absent from the Pauline corpus in the early second century.
Silence about a passage is not evidence of its absence; the
Apostolic Fathers wrote short occasional works in specific
polemical contexts, and many Pauline passages they do not cite
are nonetheless attested in P46.  P46 confirms that by c.~200~CE
both Gal~1:19 and 1~Cor~15:3--8 are textually secure, and this
paper does not argue their interpolation.

What the silence \textit{does} establish is more interesting.
Ignatius's anti-docetic polemic at c.~110~CE leverages Rom~1:3 (of
the seed of David) and Gal~4:4 (born of a woman, of a virgin)
precisely because these passages make Jesus's flesh-embodiment
load-bearing.  If Gal~1:19 had in Ignatius's understanding meant
\enquote{James the biological brother of Jesus}, it would have
served the same anti-docetic function as Rom~1:3---a
specifically family-kinship attestation of Jesus's real human
embodiment.  Ignatius does not deploy it that way.  Neither does
any earlier or contemporary Apostolic Father.  The passage is
not treated by the earliest surviving Christian writers as
evidence for Jesus's biological family.

This silence is compatible with the main paper's argument that
Gal~1:19's \textit{adelphon tou kyriou} was not read as a
biological claim in the earliest stratum of Christian reception.
It is also compatible with the alternative that Apostolic Fathers
simply had no occasion to cite it.  The data does not rule
between these readings, but it does establish that the passage's
historical-biographical weight is not attested in the extant
earliest reception.

The contrast with the citation profile of Rom~1:3---which is
demonstrably load-bearing in early anti-docetic argument---is the
empirical fact.  A historicist reading that treats Gal~1:19 as
the paradigmatic case for Jesus's historicity must explain why
the first generation of patristic writers does not cite it in
contexts where doing so would have been argumentatively useful.

\paragraph{Extension to Justin Martyr and Irenaeus.}
The finding extends to the next stratum of patristic writers.
A systematic examination of the Ante-Nicene Fathers volume~I
(covering Justin Martyr, c.~150~CE, and Irenaeus, c.~180~CE, in
addition to the Apostolic Fathers) confirms that neither Justin
nor Irenaeus cites Galatians~1:19 as \enquote{James the brother
of the Lord} nor the 1~Corinthians~15:3--8 witness list.  Justin
Martyr is known in specialist literature for his notably sparse
use of Paul generally; his arguments from scripture in the
\textit{Dialogue with Trypho} and the \textit{Apologies} rely
overwhelmingly on the Hebrew Bible and on synoptic-Gospel material
rather than on Pauline citation.  Irenaeus's \textit{Adversus
Haereses} cites Paul extensively, including Rom~1:3--4
repeatedly in its anti-Marcionite and anti-docetic argument, but
does not deploy Gal~1:19 as evidence of Jesus's biological
family.  The earliest clear citations of \enquote{James, the
brother of the Lord} as the name of a specific historical
biological sibling come from Tertullian \textit{Adversus
Marcionem} (c.~207~CE) and, more explicitly, from Jerome's
\textit{Adversus Helvidium} (c.~390~CE).  In the intervening
period between the Pauline letters (c.~50s~CE) and Tertullian,
the phrase is \textit{textually transmitted} (attested in P46) but
is \textit{not rhetorically deployed} as a biographical attestation
of Jesus's family.

(The Ante-Nicene Fathers collection does contain genealogical
discussions of James as \enquote{the Lord's cousin} or \enquote{the
Lord's half-brother} in pseudo-Ignatian material
[\textit{Letter to John Concerning the Holy Virgin Mary},
4th-c.\ forgery] and in a medieval \enquote{Fragments of Papias}
compilation that the critical apparatus identifies as a
Byzantine dictionary, not authentic Papias.  These are later
receptional material, not first- or second-century evidence.)

\paragraph{Preliminary stylometric analysis of the contested
passages.}
A first-pass stylometric analysis, reported in the project
repository (\textit{baseline/stylometry\_report.md}), offers modest
independent corroboration of the textual picture above.  The
analysis computes function-word-frequency profiles for each
Pauline letter and for individual candidate passages, then
measures the Euclidean distance of each from the global Pauline
baseline.  The method uses approximately eighty high-frequency
Koine Greek function words (conjunctions, particles, articles,
pronouns, prepositions) chosen to discriminate style rather than
content.

The method validates on its control set.  The two passages most
confidently identified as interpolations in modern scholarship---
1~Thessalonians~2:14--16 and 1~Corinthians~14:34--35---both show up
as strong stylometric outliers (distance~$\approx 0.10$ vs.\ a
per-letter baseline of~$0.015$--$0.05$).  The pre-Pauline creedal
fragments are also correctly flagged as distinct: Romans~1:3--4
(distance~$0.17$), 1~Corinthians~15:3--5 (distance~$0.16$), and
Philippians~2:5--11 (distance~$0.07$).  The rolling-window scan
additionally identifies 2~Corinthians~6:14--7:1---a passage
independently suspected of interpolation by some specialists---as
the largest outlier in the corpus.

For the historicity-relevant passages, the results are more
measured.  In wider context (Galatians~1:18--20, the three-verse
passage containing \grk{ton adelphon tou kyriou}), the
stylometric distance from Pauline baseline is $\approx 0.09$---
noticeably above baseline but \textit{less} distinctive than the
established interpolations and less distinctive than the pre-
Pauline creeds.  Galatians~4:3--5 is similarly mild
($\approx 0.11$).  1~Corinthians~9:4--6 is more distinctive
($\approx 0.18$), suggesting some stylistic movement at that
passage but not at an interpolation-diagnostic level.

The stylometric signal at Gal~1:19 alone, or 1~Cor~9:5 alone, is
not reliably distinguishable from sampling noise at the 13--19-
token scale at which individual verses can be analysed; stylometric
inference requires substantially larger windows for robust results.

The honest summary is that preliminary stylometric evidence
\textit{does not support} an interpolation hypothesis for the
principal historicity-relevant passages: those passages cluster
stylistically with their surrounding Pauline context more closely
than the established interpolations do, and meaningfully less
distinctively than the recognised pre-Pauline creedal formulae.
The interpretive argument for the community-title reading of
Gal~1:19 must therefore rely on semantic, contextual, and
corpus-linguistic grounds rather than on text-critical or
stylometric excision.

This is the honest negative result.  A proper stylometric study of
Pauline authenticity would use larger windows, multiple feature
schemes (function words, character $n$-grams, POS distributions),
and publication-standard statistical testing.  The preliminary
analysis here is offered to establish the \textit{direction} of
stylometric evidence on the contested passages; it is not an
argument against the community-title reading of Gal~1:19 as such,
but it is an argument \textit{against} defending that reading by
appeal to interpolation.

%% ------------------------------------------------------------
\subsection{The Character of Paul's Jesus}
\label{sec:paul:character}
%% ------------------------------------------------------------

The central observation, developed most thoroughly by
\citet{Doherty1999,Doherty2009}, is this: Paul's Jesus lacks the
attributes of a recently-living historical person.

Paul never names Jesus's parents, birthplace, or hometown.  He provides
no physical description.  He narrates no events from Jesus's life---no
ministry, no healings, no exorcisms, no parables, no teaching occasions,
no conflicts with authorities, no trial, no specific date or location of
execution.  He names no contemporaries who interacted with Jesus during
a putative earthly life.  He does not identify Judas as a betrayer.  He
does not name Pilate.  He does not place the crucifixion at Golgotha, in
Jerusalem, or indeed anywhere on earth.

What Paul does describe is a cosmic soteriological drama.  Jesus is a
divine figure who ``emptied himself'' and took on the ``form of a
servant'' (Philippians~2:5--11---itself likely a pre-Pauline hymn).  He
was ``declared to be Son of God in power according to the Spirit of
holiness by his resurrection from the dead'' (Romans~1:4).  He was
``handed over'' (\grk{paradidomi})---a term that in 1~Corinthians~11:23
is routinely translated ``betrayed'' to harmonize with the Gospel
narrative, but whose primary meaning is ``delivered up,'' and which in
context can refer to God handing Jesus over as a sacrificial act, not to
a human betrayal at a historical supper \citep{Doherty2009,Carrier2014}.

Paul attributes his knowledge of Jesus to two sources: revelation
(\grk{apokalypsis}---Galatians~1:12, 1:16, 2:2) and scripture
(1~Corinthians~15:3--4: ``Christ died for our sins \textit{in
accordance with the scriptures}\ldots and was raised on the third day
\textit{in accordance with the scriptures}'').  He explicitly denies
that his gospel came from any human source: ``I did not receive it from
a human source, nor was I taught it, but I received it through a
revelation of Jesus Christ'' (Galatians~1:12).

The historicist reading of Paul requires us to believe that a writer who
personally met people who had known a recently-living Jesus, and who was
engaged in active disputes about authority, practice, and doctrine
within the movement, chose to convey almost no biographical information
about that person and instead attributed his entire knowledge to
visionary experience and scriptural interpretation.  This is not
impossible, but it is precisely what we would \textit{not} expect if
Paul's Jesus were a recent historical figure, and precisely what we
\textit{would} expect if Paul's Jesus were a celestial being known
through revelation and scripture.

The pre-Pauline creedal material---Romans~1:3--4, Philippians~2:5--11,
1~Corinthians~15:3--7---deserves special attention.  These formulae are
earlier than Paul's letters themselves and thus represent the oldest
recoverable layer of Christian belief.  Historicists cite them as
evidence of early tradition about a historical person.  Examined on
their own terms, they describe a divine being who descended, took on
human likeness, died, and was exalted---a mythological pattern rather
than a biographical sketch.  However, several specific passages within
and adjacent to this creedal material pose genuine challenges for the
mythicist model that must be honestly acknowledged.

\paragraph{Galatians 4:4.}
The phrase ``born of a woman, born under the law'' (Galatians~4:4) is
cited by historicists as evidence of ordinary human birth.  The mythicist
reading---that Paul is making a theological argument about Christ
entering the domain of the Law, using language that parallels Jewish and
Hellenistic descriptions of divine beings taking on mortal form
\citep{Doherty2009}---is possible.  But it requires interpretive work.
The phrase ``born of a woman'' (\grk{genomenon ek gynaikos}) is, in
ordinary Greek, a straightforward statement about human birth.  The
mythicist must argue that Paul is using the phrase in an extraordinary
sense---as a cosmological claim about a celestial being entering the
fleshly domain---despite its ordinary denotation.  This reading is not
impossible: the context is theological (Paul is discussing the
relationship between the Law and the promise), and the passage as a
whole operates at the level of soteriology, not biography.  But the
historicist reading---that Paul simply means Jesus was born as a human
being---is the more natural one, and the mythicist case must acknowledge
this rather than dismissing the passage in a sentence.

\paragraph{Romans 1:3--4.}
Romans~1:3--4 states that Jesus was ``descended from David according to
the flesh.''  The phrase ``according to the flesh'' (\grk{kata sarka})
is doing theological work, distinguishing the domain of flesh from the
domain of spirit.  Paul uses \grk{kata sarka} elsewhere as a
cosmological category (Romans~8:4--5, 2~Corinthians~5:16).  In
2~Corinthians~5:16, Paul writes: ``even though we once knew Christ
\grk{kata sarka}, we know him thus no longer.''  If this phrase means
``as a physical human being,'' Paul is saying he once regarded Christ as
a human figure but no longer does---an extraordinary statement for a
historicist to explain and a natural one on the mythicist reading.
Davidic descent ``according to the flesh'' may describe a mythic
genealogy of a celestial being who entered the fleshly domain
\citep{Carrier2014}.  However, the historicist reading---that Paul is
making a straightforward genealogical claim---cannot be ruled out, and
the passage remains genuinely ambiguous rather than clearly favoring
either hypothesis.

\paragraph{1~Corinthians 15:3--8 and the witness list.}
The creedal formula in 1~Corinthians~15:3--8 presents the most
difficult challenge for the mythicist model within the Pauline corpus.
Paul cites a received tradition: Christ died, was buried, was raised on
the third day ``in accordance with the scriptures,'' and appeared to
Cephas, then to the Twelve, then to more than five hundred at once,
then to James, then to all the apostles, and finally to Paul himself.

Several features of this passage require careful analysis.

\textit{The verb \grk{\'ophth\=e}.}  The Greek verb translated
``appeared'' is \grk{\'ophth\=e} (the aorist passive of \grk{hora\=o},
``to see'').  This is a crucial datum.  In the Septuagint,
\grk{\'ophth\=e} is the standard term for divine theophanies---God
``appearing'' to Abraham (Genesis~12:7), to Moses at the burning bush
(Exodus~3:2), to Isaiah in the temple (Isaiah~6:1).  These are
visionary encounters with a divine being, not physical sightings of a
human body.  Paul uses this same verb for all the appearances in the
list, including his own---and Paul's own encounter with Christ was
unambiguously visionary: he describes it as a \grk{apokalypsis}
(revelation) in Galatians~1:12 and 1:16, and Acts~9 depicts it as a
blinding light and a disembodied voice.

Crucially, Paul places his own experience in the same list as the
others without any qualitative distinction: ``last of all, as to one
untimely born, \grk{\'ophth\=e} also to me'' (15:8).  The
``untimeliness'' refers to the timing, not to a difference in the
nature of the experience.  This suggests Paul understood all the
appearances---to Cephas, to the Twelve, to the five hundred---as the
same kind of event he himself experienced: a revelatory vision of the
risen Christ, not a physical encounter with a recently-living body.

\textit{The five hundred.}  The reference to ``more than five hundred
brothers at once, most of whom are still alive'' (15:6) is often cited
as evidence that Paul is appealing to verifiable witnesses of a
physical event.  This is the historicist's strongest reading.  However,
several considerations complicate it.  Collective visionary experiences
are attested in ecstatic religious movements: the Day of Pentecost
narrative in Acts~2 describes a group experience of the Spirit, and
crowd-level visionary phenomena are documented in numerous historical
and anthropological settings \citep{Lewis1971}.  Paul may be reporting
a communal ecstatic event---a mass visionary experience that the
community interpreted as an appearance of the risen Christ.

Additionally, \citet{Carrier2014} has argued that the witness list as
we have it may represent an expansion of an originally shorter creed.
The core formula (15:3--5: ``died\ldots buried\ldots raised\ldots
appeared to Cephas, then to the Twelve'') is metrically and
structurally distinct from the expansions that follow (the five hundred,
James, ``all the apostles,'' and Paul himself).  If the list grew over
time as the tradition developed, the ``five hundred'' may be a later
addition rather than part of the original creed Paul received.  This is
speculative, but the textual seams in the passage are real.

\textit{Assessment.}  The 1~Corinthians~15 witness list is the single
strongest piece of evidence for the historicist case within the Pauline
corpus.  The specificity of named individuals (Cephas, James), a
defined group (the Twelve), and a large crowd is difficult to explain
under pure mythicism.  The mythicist reading---that these are all
visionary experiences, as Paul's own use of \grk{\'ophth\=e} and his
inclusion of his own visionary encounter in the same list
suggest---is genuinely plausible, but the appeal to living witnesses
(``most of whom are still alive'') introduces a historical-verification
claim that resists easy dismissal.  This passage is a genuine unresolved
problem for the mythicist hypothesis, and intellectual honesty requires
saying so.  It does not on its own establish historicity---visionary
movements can generate elaborate witness traditions---but it is the
point where the historicist case is most difficult to answer.

%% ------------------------------------------------------------
\subsection{The Pastoral Silence}
\label{sec:paul:silence}
%% ------------------------------------------------------------

The strongest form of the argument from the Pauline corpus is not the
absence of biographical detail \textit{per se}---one might excuse this
as Paul's theological focus---but the absence of biographical detail
\textit{in contexts where it would be decisive}.

Paul addresses practical questions of community life throughout his
letters.  In 1~Corinthians alone, he takes up marriage and divorce
(ch.~7), food offered to idols (chs.~8--10), proper conduct at
communal meals (ch.~11), the exercise of spiritual gifts (chs.~12--14),
and the resurrection of the dead (ch.~15).  In Romans, he addresses
relations with governing authorities (ch.~13), disputes about dietary
observance (ch.~14), and the status of Israel (chs.~9--11).  In
Galatians, the question is whether Gentile converts must follow Jewish
law.

In \textit{every one} of these disputes, a single authoritative saying
of the historical Jesus would be the trump card.  A teacher who had
recently lived and spoken on matters of religious practice would
presumably have left sayings relevant to these exact questions---and if
Paul knew such sayings, citing them would instantly settle the argument.

Paul almost never does this.  In the entirety of the undisputed letters,
there are only two passages where Paul appears to cite a ``command of the
Lord'' on a practical matter: 1~Corinthians~7:10 (on divorce) and
1~Corinthians~9:14 (on apostolic support).  These two passages require
careful treatment because they constitute a genuine difficulty for the
mythicist model.

In 1~Corinthians~7:10--11, Paul writes: ``To the married I give this
command---not I but the Lord---that the wife should not separate from
her husband.''  He distinguishes this from his own opinion (7:12: ``To
the rest I say---I and not the Lord\ldots''), indicating that he regards
the divorce command as carrying a different and higher authority than
his personal judgment.

In 1~Corinthians~9:14, Paul writes: ``In the same way, the Lord
commanded that those who proclaim the gospel should get their living by
the gospel.''

The mythicist reading treats these as prophetic oracles---commands
received from the risen, cosmic Lord through spiritual revelation, not
transmitted sayings of a historical teacher \citep{Aune1983}.  Early
Christian prophets did issue pronouncements in the name of the Lord
(Revelation is full of such pronouncements: ``the Spirit says to the
churches\ldots''), and Paul understood himself to be in direct
communication with the risen Christ (Galatians~1:12, 2:2;
2~Corinthians~12:1--4).

This reading is \textit{possible}, but it faces a genuine problem.  If
all dominical commands were prophetic oracles available on demand through
the Spirit, Paul would have an unlimited supply of them---and we would
expect him to invoke ``commands of the Lord'' far more frequently than
twice across seven letters.  The scarcity is more naturally explained if
Paul had access to a small number of transmitted sayings attributed to
a founder, two of which he found applicable, rather than an open channel
of prophetic revelation he invoked only twice.

On the other hand, the historicist reading also faces a difficulty: if
Paul had access to a body of transmitted historical sayings, the two
citations amid dozens of missed opportunities is strikingly low.  A
writer with a catechism of Jesus sayings would cite them far more
often.  The pattern---two and only two---is anomalous on either
hypothesis, but it is \textit{more} anomalous on the historicist reading
(which predicts many citations) than on the mythicist reading (which
predicts very few, arising from prophetic oracles that happened to match
later Gospel sayings or from a tiny core of early community rulings
attributed to the Lord).

The honest assessment is that 1~Corinthians~7:10 and 9:14 resist easy
explanation under either model.  They are the strongest evidence within
the pastoral silence analysis for some form of transmitted dominical
tradition---but the tradition is so sparse that it does not clearly
indicate a recently-living teacher with an extensive body of teaching.

The silence is most glaring in cases where we know, from the later
Gospel tradition, that Jesus supposedly spoke directly to the issue at
hand.  Paul discusses food laws without citing any saying like Mark~7:15
(``there is nothing outside a person that by going in can defile'').  He
discusses relations with Rome without citing anything like ``render unto
Caesar.''  He discusses the Law at extraordinary length in Galatians and
Romans without citing Jesus's supposed pronouncements on the Law from
the Sermon on the Mount.

The historicist explanation---that Paul's letters are occasional
correspondence, not systematic theology, and we should not expect
exhaustive citation of dominical sayings---has a limit.  Paul is not
merely writing chatty letters.  He is \textit{arguing}, often
polemically, for specific positions against specific opponents.  He is
asserting authority.  In such contexts, the failure to cite the
authoritative words of the founder, when those words would be
dispositive, demands explanation.

A more sophisticated historicist response deserves acknowledgment: Paul
may have deliberately avoided citing transmitted sayings because his
authority rested on direct revelation from the risen Christ, not on
second-hand memory of a historical teacher.  Citing earthly sayings
would have subordinated his apostolic authority to that of the Jerusalem
apostles who had known Jesus personally---precisely the hierarchy Paul
resists throughout Galatians.  On this reading, Paul's silence is
strategic, not evidential.

This explanation is \textit{ad hoc}, but no more so than the mythicist
reading of Galatians~1:19 as a community title.  However, it faces the
difficulty discussed above: Paul \textit{does} cite ``commands of the
Lord'' in two cases (1~Corinthians~7:10 and 9:14), undermining the
claim that his silence is a deliberate strategy.  The
pattern---two citations amid dozens of missed opportunities---is more
naturally explained by a writer who has very few such sayings available
than by a writer who has many but strategically withholds them, though
neither reading is without difficulty (see the detailed analysis above).

\citet{Doherty2009} summarizes the mythicist reading with characteristic
directness: Paul's letters read as the correspondence of a man who has
never heard of a human Jesus.  Every element of his Christology can be
derived from Jewish scriptural exegesis, Hellenistic soteriology, and
visionary experience.  This remains the most natural reading of the
pastoral silence, even if it is not the only possible one.

%% ------------------------------------------------------------
\subsection{\grk{Adelphos tou Kyriou}: ``Brother of the Lord''}
\label{sec:paul:brother}
%% ------------------------------------------------------------

If the foregoing analysis is correct, the strongest historicist evidence
from the Pauline corpus---beyond the genuinely challenging passages
discussed above (Galatians~4:4, Romans~1:3, 1~Corinthians~15:3--8)---is
a single passage: Galatians~1:19.

Paul, describing his first visit to Jerusalem, writes: ``I did not see
any other apostle except James, \grk{ton adelphon tou kyriou}---the
brother of the Lord.''  The historicist reading is straightforward: James
was the biological brother of Jesus; if Jesus had a biological brother,
Jesus was a biological person.  \citet{Ehrman2012} considers this
passage virtually decisive on its own.

The passage deserves the most careful scrutiny, precisely because it is
being asked to carry the entire weight of the historicist case.

\paragraph{The semantic range of \grk{adelphos}.}
Paul uses \grk{adelphos} (brother) and its plural \grk{adelphoi}
(brothers) extensively throughout his letters as a standard term for
fellow members of the Christian community.  This is not disputed.
Romans alone contains the term dozens of times in this sense.  The word
does not inherently denote biological kinship in Pauline usage; its
default meaning in the Pauline corpus is ``fellow believer.''

This reading is not a novel mythicist inference.  It is the standard
lexicographical treatment.  The \citet{CambridgeGreekLexicon2021}---
the first major English-language replacement for the Intermediate
Liddell-Scott in over a century, covering authors from Homer through
the early second century~CE---distinguishes two principal senses of
\grk{adelphos}: (1) \enquote{male child from the same womb, brother,}
cited from Homer onward; and (2) a figurative or hyperbolic use
subdivided into \enquote{brother} as a person regarded with especial
affection (Xenophon, NT), \enquote{(ref.\ to a fellow Christian) NT,}
and \enquote{(gener., ref.\ to an associate or neighbour) NT.}  The
community-of-believers sense is therefore not an interpretive
stretch but a separately-listed, lexicographically-recognized sense
attested specifically in the New Testament corpus within which
Paul's letters circulate.  The question for Galatians~1:19 is which
of these Cambridge-Lexicon-recognized senses Paul is deploying
there---sense 1, or sense 2b---and the argument that follows in
this section contends that sense 2b is the more parsimonious
reading given Paul's own corpus-wide distribution and the Greek
religious-associative precedent.

A further lexicographical observation worth recording is the
pattern of inclusions and omissions within the Cambridge Lexicon
for the \grk{adelph-} family of compounds Paul and other NT
authors use.  Examination of the three relevant pages---page~18
(the \grk{adelph-} section), page~1462 (\grk{phil-} compounds),
and page~1518 (\grk{pseud-} compounds)---yields a strikingly clean
division.

Cambridge includes every \textit{classical} and
\textit{pre-NT} \grk{adelph-} word: the parent noun \grk{adelphos}
(sense 1 biological, sense 2 figurative subdivided to include a
specific \enquote{(ref.\ to a fellow Christian) NT} sense); the
feminine \grk{adelph\=e} (with a softer NT extension for \enquote
{a woman regarded w.\ especial affection}); the adjective
\grk{adelphos} \grk{h\=e} \grk{on} (\enquote{kindred, twin, akin});
the derived adjective \grk{adelphikos} (\enquote{brotherly}); the
compound \grk{phil-adelphos} (\enquote{brother-loving,} cited only
from Xenophon, Plutarch, Sophocles, and Polybius, no NT citation);
the kinship nouns \grk{adelphid\=e} (niece) and \grk{adelphidous}
(nephew); the diminutive \grk{adelphidion} (\enquote{little
brother}); the compound \grk{adelpheoktonos} (fratricide); and the
verb \grk{adelphiz\=o} (\enquote{to call someone `brother' in order
to ingratiate oneself}).  Cambridge also includes on the same
pages NT-specific compound nouns with non-\grk{adelph-} heads, such
as \grk{pseudo-proph\=et\=es} (\enquote{false prophet,} NT-only) and
\grk{pseudo-martys} (\enquote{false witness,} Pl.\ and NT).

Cambridge does \textit{not} include, as separate headwords:
\grk{philadelphia} (\enquote{brotherly love,} NT-dominant---Paul
at Rom~12:10 and 1~Thess~4:9, plus 1~Pet~1:22, 2~Pet~1:7, Heb~13:1);
\grk{pseudadelphos} (\enquote{false brother,} Paul at Gal~2:4 and
2~Cor~11:26, otherwise unattested pre-Christian); and
\grk{adelphot\=es} (\enquote{brotherhood,} 1~Pet~2:17 and 5:9,
otherwise unattested pre-Christian).  All three omissions are
lexicographical judgements, not editorial policy---Cambridge
demonstrates on the same pages that it is willing to list
NT-specific compounds and parallel abstract nouns (e.g.,
\grk{philandria} on the same page as the absent
\grk{philadelphia}).

The omissions are lexicographical judgements, not oversights.  They
fit the attestation profile of the two omitted words: both
\grk{philadelphia} and \grk{pseudadelphos} are essentially
first-century Christian technical terms.  Paul uses
\grk{philadelphia} at Romans~12:10 and 1~Thessalonians~4:9; the
other principal NT occurrences are 1~Peter~1:22, 2~Peter~1:7, and
Hebrews~13:1; the pre-Christian attestations (Plutarch's
\textit{Peri~Philadelphias}, a few others) are rare enough that
Cambridge's editors did not judge them sufficient for a separate
lemma.  Paul uses \grk{pseudadelphos} at Galatians~2:4 and
2~Corinthians~11:26; the pre-Christian attestations are effectively
nil.  The pattern is consistent with the paper's broader
observation that the community-of-brothers vocabulary was
specialized in early-Christian discourse from the pre-existing
Greek raw materials of classical kinship language.  The
\textit{adjective} \grk{phil-adelphos} predates Christianity and
has a robust classical pedigree; the \textit{abstract noun}
\grk{philadelphia} and the \textit{compound noun}
\grk{pseudadelphos} are Christian-era developments that extended
this pre-existing vocabulary.

The evidential force of this observation for §\ref{sec:paul:brother}
is modest but real: Paul writes within, and contributes to, a
specifically religious-community dialect of Greek that built up a
family of \grk{adelph-} compound terms around the community-sense
of the parent noun.  This is the linguistic signature of a
religious movement using (and extending) the available associative-
kinship vocabulary of its ambient Greek---exactly the profile
§\ref{sec:paul:brother} argues that \grk{adelphos tou kyriou} fits.

The historicist response is that ``brother of the Lord'' (\grk{adelphos
tou kyriou}) is a \textit{specific construction} that adds the genitive
\grk{tou kyriou}, distinguishing it from the ordinary fraternal usage.
This is said to mark it as a kinship claim rather than a community
designation.  But this argument assumes what it needs to demonstrate.
The addition of ``of the Lord'' is equally consistent with a
\textit{title} distinguishing a particular subgroup or rank within the
movement from ordinary \grk{adelphoi}---``Brothers of the Lord'' as a
named class, analogous to ``apostles'' or ``pillars'' (Galatians~2:9).

\paragraph{1~Corinthians 9:5 and the plural ``brothers of the Lord.''}
The decisive test case is 1~Corinthians~9:5, where Paul writes: ``Do we
not have the right to be accompanied by a believing wife, as do the
other apostles and the brothers of the Lord [\grk{hoi adelphoi tou
kyriou}] and Cephas?''

Paul is listing categories of authoritative figures: (1)~the other
apostles, (2)~the brothers of the Lord, and (3)~Cephas.  Three distinct
groups or individuals, cited in parallel as recognized authorities
within the movement.  ``The brothers of the Lord'' is \textit{plural}
and functions as a class label.

If this phrase means ``biological siblings of Jesus,'' then there exists
an entire group of biological brothers, all sufficiently prominent to be
cited as a recognized authority class alongside the apostles and Peter
individually.  The Gospels eventually supply names for such
brothers---James, Joses, Judas, Simon (Mark~6:3)---but this is in
texts written decades later, and the question is whether the Gospel
detail was generated \textit{from} the community title (once
``brothers of the Lord'' had been reinterpreted as biological kinship
after historicization of the Jesus figure) or whether the community
title reflects pre-existing biographical knowledge.

The simpler reading is available: ``Brothers of the Lord'' is a title
for a specific group or rank within the early Christian movement---an
inner circle, a particular faction, or a designated class---distinguished
from ``apostles'' and from individual leaders like Cephas.  Paul
identifies ``James the brother of the Lord'' in Galatians~1:19 not by
biological relationship but by group affiliation: James of the Brothers
of the Lord, as one might say ``James of the Jerusalem council'' or
``James the Elder.''

\paragraph{``The Lord'' as a title, not a personal name.}
In Pauline usage, ``the Lord'' (\grk{ho kyrios}) is not a name for a
historical person who recently walked in Galilee.  It is a title for the
risen, cosmic, divine Christ---the figure Paul encountered in
revelation, who sits at the right hand of God, who will return at the
\grk{parousia}.  Every other Pauline use of ``the Lord'' refers to this
exalted, divine figure.  ``Brother of the Lord'' therefore means
``brother of the cosmic Christ,'' which is far more naturally a
community designation---marking special proximity or authority relative
to the divine patron---than a biological kinship claim about a recent
human being.

\paragraph{The interpolation possibility.}
Galatians contains recognized interpolation issues.  The clause
\grk{ton adelphon tou kyriou} could be a marginal gloss---a later
scribe's clarifying note identifying which James Paul meant, written in
the margin and subsequently incorporated into the text.  This is
precisely the kind of editorial intervention that occurred routinely in
the manuscript tradition, as the pseudepigraphical letters and the
known interpolations in the authentic letters demonstrate.  By the time
later tradition had firmly identified James as the biological brother
of Jesus (a tradition well established by the time of Hegesippus in the
mid-second century), a scribe copying Galatians would have natural
motivation to add the clarifying tag.  This possibility cannot be
confirmed, but neither can it be excluded, and it means we cannot be
certain the phrase is authentically Pauline.

\paragraph{Assessment.}
The biological reading of Galatians~1:19 is \textit{a} possible
reading.  It is not the \textit{only} possible reading, and when the
full Pauline context is considered---the routine use of \grk{adelphos}
for fellow believers (117 community-sense occurrences across the seven
undisputed letters), the plural ``brothers of the Lord'' functioning as
a group label in 1~Corinthians~9:5, the cosmic rather than biographical
referent of ``the Lord'' throughout Paul, and the possibility of
interpolation---it is the \textit{least parsimonious} reading.  The
community-title reading requires no special pleading; the biological
reading requires that Paul used \grk{adelphos} in a unique sense in this
one passage, that ``the Lord'' refers to a historical person in a way it
does nowhere else in Paul, and that the plural usage in
1~Corinthians~9:5 (which treats ``brothers of the Lord'' as a group
parallel to ``apostles'') is coincidentally misleading.

\paragraph{The absence of a biological-sibling register in Paul.}
A stronger point follows from a complete lexical survey.  Across the
seven undisputed Pauline letters there are \textit{one hundred
twenty-one} occurrences of the \grk{adelph-} family of lexemes
(\grk{adelphos}, \grk{adelph\=e}, \grk{pseudadelphos},
\grk{philadelphia}).  Of these, \textit{zero} unambiguously denote a
biological sibling outside the two disputed \grk{tou kyriou} passages.
Every other occurrence either addresses the community (``brothers, I
urge you\ldots''), names a coworker (``Timothy the brother,'' ``Apollos
the brother,'' ``Titus our brother''), identifies a believing spouse
(1~Corinthians~7:12--15), takes an explicit modifier such as
\grk{kata sarka} for ethnic kinship (Romans~9:3) or \grk{pseudo-} for
false brothers (2~Corinthians~11:26, Galatians~2:4), or is otherwise
irreducibly ambiguous.  Paul has no biological-sibling register of
\grk{adelphos} from which a kinship reading of Galatians~1:19 could be
drawn.  A complete enumeration of every occurrence, colour-coded by
classification and embedded in the running SBLGNT text, is provided in
the accompanying supplementary material.

\paragraph{Comparison with contemporary Greek usage.}
A supplementary study establishes the baseline rate of community-sense
\grk{adelphos} in contemporary non-Christian Greek by random sampling
across six corpora spanning documentary papyri (100 BCE--70 CE), the
Septuagint, Philo of Alexandria, Josephus, and pagan literary authors
both pre-Pauline (Polybius, Diodorus, Strabo) and contemporary
(Plutarch, Epictetus, Dio Chrysostom).  Across 405 randomly-sampled
tokens (three literary buckets at $n=100$, three at $n=35$ pending
external re-sampling), the community/voluntary-association rate (A)
is \textit{0.7\%} with 95\%~Wilson confidence interval $[0.3\%,~2.2\%]$.
The comparable Pauline rate is \textit{96.7\%} ($[91.8\%,~98.7\%]$).
The two confidence intervals do not overlap.

This sampled baseline, however, \textit{excludes Greek inscriptions},
and as \citet{Harland2003,Harland2005} has documented, inscriptions are
precisely where community-sense \grk{adelphoi} concentrates.  Voluntary
cult-associations (\grk{thiasoi}), mystery-cult initiate groups,
professional guilds, and some civic brotherhoods routinely designated
their fellow members as \grk{adelphoi} in stone inscriptions across
the Greek East.  The Kloppenborg-Ascough corpus of Greco-Roman
associations \citep{KloppenborgAscough2011} catalogs hundreds of
relevant texts.  Factoring the inscriptional register into the baseline
would raise the non-Christian community-sense rate from roughly 1\%
toward something in the range of 5--15\%, and the Pauline-to-baseline
ratio would correspondingly contract from the nominal $\sim$130$\times$
toward the 10--20$\times$ range.  That contraction is material, and
any responsible reading of the baseline study should apply it.  The
main substantive conclusion---that Paul's community-sense usage is
distinctive against his ambient register---survives.  The quantitative
shape of the distinction does not.

\paragraph{The precedent for community-title \grk{adelphoi}.}
Harland's work changes the argumentative frame in a more interesting
way than the numerical correction alone suggests.  If Greek voluntary
associations routinely used \grk{adelphoi} as a community-title---a
designation of fellow members of an elective religious or cultic
group---then the reading of \grk{hoi adelphoi tou kyriou} in
1~Corinthians~9:5 and \grk{ton adelphon tou kyriou} in Galatians~1:19
as community titles requires no novel semantic stretch.  It instantiates
a well-attested Greek pattern.  Harland's \citeyear{Harland2005} article
discusses specifically inscriptional uses of \grk{adelphoi tou theou}
(\enquote{brothers of the god}) and similar constructions designating
members of a cult association around a patron deity.  The Pauline
construction---\enquote{brothers of the Lord}---is formally identical:
fellow members of a religious community whose focal figure is the
cosmic Christ whom Paul calls \grk{ho kyrios}.  Far from being an
unusual or suspicious move, a community-title reading of Gal~1:19 is
the reading that places Paul most squarely within his contemporary
religious-associative milieu.

Both points combine: Paul's usage is distinctive against everyday
Greek (where biological meaning dominates), but fits naturally within
the well-attested Greek religious-association register (where
community-title \grk{adelphoi} is normal).  Gal~1:19 as a
community-title is not a forced reading; it is the culturally default
reading for the specific subregister in which Paul is writing.
Details of the baseline study, the coding scheme, its limitations, and
the Harland-sensitivity analysis are given in the companion document
\textit{Corpus-Linguistic Baseline for Adelphos in Contemporary Greek}.

This is a stronger observation than the frequency argument above.  The
historicist response to a frequency claim is that \grk{adelphos tou
kyriou} is a \textit{specifically modified} construction, and the 117
generic uses of plain \grk{adelphos} do not govern the kinship reading
of the genitive phrase.  But that response assumes the existence of a
separate, well-attested biological-sibling usage from which the
modified form is drawn.  In Paul, there is no such usage.  The
biological reading of Galatians~1:19 therefore requires positing a
lexical register that Paul exhibits nowhere else in his authentic
corpus---a hidden kinship sense invoked once, in a single phrase,
against a backdrop of 117 community uses and zero biological ones.

Galatians~1:19 is compatible with historicity; it does not require it.
And a single phrase, read against a corpus in which biological sibling
is a register the author never uses elsewhere, embedded in a letter
collection that otherwise describes a celestial revelatory figure with
no biographical specifics, is not the foundation on which the most
consequential historical claim in Western civilization should rest.


%% ============================================================
\section{The Gospel Sources}
\label{sec:gospels}
%% ============================================================

The Gospels---Mark, Matthew, Luke, and John---are the primary narrative
sources for the life of Jesus.  They are routinely treated as evidence
for historicity, both by the consensus and in popular understanding.
The question is whether this treatment is warranted.

It is uncontroversial in mainstream scholarship that the Gospels are not
eyewitness accounts \citep{Ehrman2009}.  They are anonymous works
(the names Mark, Matthew, Luke, and John are second-century
attributions), written in Greek by educated authors, forty to seventy
years after the events they purport to describe, in communities far
removed from rural Galilee.  They are theological documents with
literary and rhetorical purposes.  The question is whether, beneath the
theology and literary construction, a historical core can be recovered.

The argument of this section is that it cannot---not because the Gospels
are ``unreliable'' in the colloquial sense, but because their narrative
content is demonstrably constructed from identifiable prior
sources---Hebrew Bible passages, Hellenistic literary models, and each
other---in a way that leaves no residue requiring a historical event to
explain.

%% ------------------------------------------------------------
\subsection{Mark as Literary Composition}
\label{sec:gospels:mark}
%% ------------------------------------------------------------

The Gospel of Mark, dated to approximately 70~CE, is the earliest
surviving narrative of Jesus's life and is the primary source for
Matthew and Luke, who independently copy and revise large portions of
it.  Mark is therefore the foundation of the entire Synoptic tradition.
If Mark is a literary composition rather than a record of historical
memory, the downstream Gospels inherit that character.

That Mark is a sophisticated literary construction is increasingly
recognized.  \citet{MacDonald2000} argues in \textit{The Homeric Epics
and the Gospel of Mark} that Mark consciously models episodes on Homeric
narratives---the calming of the storm on the Sea of Galilee on
Odysseus's sea adventures, the feeding of the five thousand on
Circe's feast, the young man who flees naked at Gethsemane
(Mark~14:51--52) on Elpenor's fall.  MacDonald's mimesis thesis remains
debated in its specifics, but the broader principle it
illustrates---that Mark is composing narrative from literary templates,
not recording eyewitness testimony---is far less controversial.

The Hebrew Bible templates are more important and less disputed.  Mark's
Jesus is modeled extensively on the Elijah--Elisha cycle from 1--2~Kings:

\begin{itemize}[leftmargin=2em]
\item Elijah calls Elisha from his work; Jesus calls disciples from
their fishing (Mark~1:16--20; \cf\ 1~Kings~19:19--21).
\item Elisha multiplies loaves to feed a crowd with leftovers
(2~Kings~4:42--44); Jesus feeds five thousand with leftovers
(Mark~6:35--44).
\item Elisha raises the dead (2~Kings~4:32--37); Jesus raises the dead
(Mark~5:35--43).
\item Elisha heals a leper (2~Kings~5:1--14); Jesus heals a leper
(Mark~1:40--45).
\item Elijah parts the water and crosses on dry ground
(2~Kings~2:8); Jesus walks on water (Mark~6:45--52).
\end{itemize}

The pattern is not one of incidental similarity; it is sustained
typological recapitulation, in which the new prophet systematically
surpasses the old.  This is a known literary technique in Jewish
writing---the presentation of a new figure as the antitype of an older
scriptural figure---and it requires no historical events to explain.  It
requires a literate author with a copy of the Septuagint and a
theological purpose.

Mark's narrative structure also reveals compositional artifice.  The
Gospel is organized around a journey from Galilee to Jerusalem,
punctuated by three passion predictions (8:31, 9:31, 10:33--34), each
followed by the disciples' misunderstanding and Jesus's corrective
teaching.  This triadic structure is literary, not mnemonic.  The
``messianic secret'' motif---Jesus repeatedly commands silence about
his identity---is a narrative device managing the tension between
Mark's theology (Jesus is the Messiah) and historical plausibility
(if Jesus publicly claimed messiahship, everyone would know)
\citep{Wrede1901}.  \citet{Wrede1901} demonstrated this over a century
ago; it remains one of the most important insights in Gospel studies,
and it points directly to Mark as a literary architect, not a compiler
of tradition.

%% ------------------------------------------------------------
\subsection{The Q Hypothesis and Its Implications}
\label{sec:gospels:q}
%% ------------------------------------------------------------

Matthew and Luke share approximately 235 verses of material not found
in Mark, much of it sayings attributed to Jesus.  The dominant
explanation in scholarship is the Two-Source Hypothesis: Matthew and
Luke independently used Mark and a second source, designated Q (from
German \textit{Quelle}, ``source''), consisting primarily of sayings
\citep{Kloppenborg1987}.

Q's existence is an inference, not a fact.  No manuscript of Q has ever
been found.  \citet{Goodacre2002} and others have argued that the
Q material can be explained by Luke's direct use of Matthew, eliminating
the need for a hypothetical source.  The Farrer hypothesis, as this
alternative is known, has gained adherents but remains a minority
position.

For the historicity question, the Q debate has different implications
depending on its resolution, but both outcomes create difficulties for
the standard historicist case---though in different ways.

If Q existed, it was a sayings collection with no passion narrative, no
birth narrative, no death-and-resurrection account, and no biographical
framework.  The earliest recoverable layer of Q (designated Q$^1$ by
\citealt{Kloppenborg1987}) consists of wisdom sayings and
instructions---the kind of material that can be attributed to a
philosophical tradition or a personified Wisdom figure without requiring
a specific historical founder.  That the earliest sayings layer lacks
any narrative of the figure's life or death is striking.  It is
consistent with a tradition that originated as wisdom instruction
attributed to a divine or personified source and was only later embedded
in a biographical narrative.  On this scenario, the historicist retains
the ``multiple independent sources'' argument (Mark plus Q as separate
streams), but the character of the earliest recoverable source points
away from a biographical tradition and toward a wisdom-school or
revelatory origin.

If Q did not exist and Luke used Matthew directly (the Farrer
hypothesis), then the ``multiple independent sources'' argument
collapses at a different point: the supposed independence of Matthew and
Luke disappears, and both are reworkings of Mark.  On this scenario,
the non-Markan material they share reflects literary dependence, not
independent access to historical tradition.  Ehrman and others
\citep{Ehrman2012} cite source plurality as a cornerstone of
historicity, but under the Farrer hypothesis, that plurality is
substantially reduced.

The two scenarios are not interchangeable.  If Q existed, the mythicist
draws support from the \textit{character} of the early tradition (wisdom
sayings, no biography).  If Q did not exist, the mythicist draws
support from the \textit{collapse of independence} among the sources.
The paper does not require a commitment to either scenario; it notes
that each, for different reasons, undermines a different pillar of the
historicist case.

%% ------------------------------------------------------------
\subsection{Matthew and Luke: Independence or Dependence?}
\label{sec:gospels:synoptic}
%% ------------------------------------------------------------

The treatment of Matthew and Luke as independent witnesses is central
to the historicist ``multiple attestation'' argument.  But even on the
Two-Source Hypothesis, their independence is limited.  Both copy Mark
extensively---Matthew reproduces approximately 90\% of Mark's content,
Luke approximately 65\%.  Their agreement with each other, beyond what
they share with Mark, is attributed to Q.  Their unique material (M and
L) may represent independent traditions---or it may represent the
evangelists' own compositions.

The birth narratives illustrate the problem acutely.  Matthew and Luke
both include birth narratives, but they are mutually contradictory:

\begin{itemize}[leftmargin=2em]
\item Matthew has the family residing in Bethlehem, fleeing to Egypt,
and relocating to Nazareth (Matthew~2:1--23).  Luke has the family
residing in Nazareth, traveling to Bethlehem for a census, and
returning to Nazareth (Luke~2:1--39).
\item Matthew provides a genealogy tracing Jesus through Solomon
(Matthew~1:1--17); Luke provides a different genealogy tracing Jesus
through Nathan (Luke~3:23--38).  The two genealogies diverge
completely after David.
\item Matthew includes the Magi, the star, and Herod's massacre of
infants.  Luke includes shepherds, angels, and a manger.  Neither
mentions any element found in the other.
\end{itemize}

These are not minor discrepancies reconcilable by harmonization.  They
are fundamentally incompatible accounts constructed to serve different
theological purposes: Matthew's narrative fulfills proof-texts from
Hosea~11:1 (``out of Egypt I called my son''), Jeremiah~31:15
(Rachel weeping for her children), and Micah~5:2 (the ruler from
Bethlehem); Luke's narrative places Jesus's birth in the context of
Roman imperial history and presents him as savior of the world in
contrast to Augustus.  Both are theological compositions, not
historical reports, and their mutual contradiction demonstrates that
the authors were not working from shared historical memory.

The same pattern---theological composition from scriptural
templates---pervades the unique material in both Gospels.  Matthew's
Sermon on the Mount is structured as a new Torah delivered from a new
Sinai by a new Moses.  Luke's parables (the Good Samaritan, the
Prodigal Son) are sophisticated literary creations with no marks of
oral transmission from a historical figure.

%% ------------------------------------------------------------
\subsection{John and the Johannine Tradition}
\label{sec:gospels:john}
%% ------------------------------------------------------------

The Gospel of John, dated to approximately 90--100~CE, is routinely
acknowledged to be the least historically reliable of the four Gospels
\citep{Brown1966}.  Its Jesus speaks in long theological discourses
entirely different in style and content from the Synoptic sayings.  Its
chronology diverges from the Synoptics (a three-year ministry rather
than one; the crucifixion on the day of preparation for Passover rather
than Passover itself).  Its theological framework---the pre-existent
Logos who descends, reveals, and returns---is explicitly mythological.

John is sometimes cited for corroborating details that differ from the
Synoptics, on the theory that independent traditions converge.  But
John's relationship to the Synoptics is disputed: many scholars see
evidence that John knew Mark or at least Markan tradition
\citep{Bauckham1998}, which compromises claims of independence.  Even
where John appears independent, the ``corroboration'' consists of
different theological constructions, not convergent historical testimony.

For the historicity question, John adds nothing.  It is the most overtly
mythological Gospel, the most distant from any putative historical
events, and the least useful as evidence for anything other than the
state of Christological reflection at the end of the first century.

%% ------------------------------------------------------------
\subsection{Scriptural Construction of the Passion Narrative}
\label{sec:gospels:passion}
%% ------------------------------------------------------------

The passion narrative---the account of Jesus's arrest, trial,
crucifixion, and burial---is often treated as the most historically
grounded portion of the Gospel tradition.  \citet{Crossan1995} and
others have argued for a pre-Markan passion narrative that may preserve
early tradition.  The historicity of the crucifixion itself is
considered one of the two ``almost universally accepted'' facts about
Jesus (the other being the baptism by John) \citep{Ehrman2012}.

But examination of the passion narrative reveals that it is constructed,
point by point, from Hebrew Bible texts.  The following correspondences
are well established in scholarship:

\begin{itemize}[leftmargin=2em]
\item \textbf{Psalm~22:} ``My God, my God, why have you forsaken me?''
(Mark~15:34 / Psalm~22:1).  ``They divide my garments among them, and
for my clothing they cast lots'' (Mark~15:24 / Psalm~22:18).
``All who see me mock me; they hurl insults, shaking their heads''
(Mark~15:29 / Psalm~22:7).  ``They pierce my hands and my feet''
(Psalm~22:16, reflected in the crucifixion itself).

\item \textbf{Isaiah~53:} ``He was oppressed and afflicted, yet he did
not open his mouth'' (Isaiah~53:7; \cf\ Jesus's silence before
Pilate, Mark~15:5).  ``He was numbered with the transgressors''
(Isaiah~53:12; \cf\ Mark~15:27, crucified between two bandits).
``He bore the sin of many'' (Isaiah~53:12; \cf\ the soteriological
interpretation of the crucifixion throughout the New Testament).

\item \textbf{Wisdom of Solomon 2--5:} The righteous one is
persecuted by the wicked, who test him with insult and torture; he
claims God as his father; the wicked condemn him to a shameful death;
after death he is vindicated and revealed as God's son.  The
structural parallel with the passion and resurrection narrative is
comprehensive \citep{Nickelsburg1972}.

\item \textbf{Zechariah:} The entry on a donkey (Zechariah~9:9 /
Mark~11:1--7).  Thirty pieces of silver (Zechariah~11:12--13 /
Matthew~26:15).  ``They will look on the one they have pierced''
(Zechariah~12:10 / John~19:37).  The shepherd struck, the sheep
scattered (Zechariah~13:7 / Mark~14:27).

\item \textbf{Daniel~7:} The Son of Man coming on clouds with power
(Daniel~7:13--14 / Mark~14:62, Jesus's declaration before the
Sanhedrin).

\item \textbf{Amos~8:9:} ``I will make the sun go down at noon and
darken the earth in broad daylight'' (\cf\ the darkness over the land
at the crucifixion, Mark~15:33).
\end{itemize}

The question is directional: did historical events occur and then get
described in scriptural language, or did scriptural passages generate
the narrative?  The historicist assumption is the former.  But the
density of correspondence---virtually every detail of the passion can
be traced to a specific scriptural source---makes the compositional
explanation more parsimonious.  The narrative reads as a midrashic
construction: a story built from scripture to demonstrate that the death
of the Christ figure fulfilled prophetic expectation.  This is a
well-attested genre in Jewish literary practice
\citep{Neusner1987,Vermes1973}.

\citet{Crossan1995} himself, despite defending historicity, acknowledged
that the passion narrative may be ``prophecy historicized'' rather than
``history remembered''---a remarkably candid concession from within the
historicist camp.  The distinction is precisely the one at issue: if the
passion is prophecy historicized, then its narrative content is
generated from scripture, not from historical memory, and it provides no
independent evidence for historical events.

\paragraph{What scriptural construction establishes, and what it does not.}
It is important to be precise about the logical step the argument is
taking, because a historicist critic such as \citet{Allison2010} would
correctly insist that scriptural templating of a narrative does not
entail the non-existence of the subject.  Historical figures have been
described, even in antiquity, through scriptural and typological
grammars.  Jewish historiography routinely employs Hebrew Bible
typology to narrate the lives of real people (Maccabean heroes as new
Judges, exilic returnees as new Exodus).  Pagan biographical narrative
likewise uses conventional \textit{theios an\=er} templates to narrate
figures whose historical existence is not in question.  That the
Gospel passion is constructed from Psalm~22 and Isaiah~53 therefore
does not on its own prove that no crucifixion-of-Jesus occurred.

The conclusion the argument actually supports is narrower but
sufficient.  It is that the Gospel passion narrative is not
\textit{independent historical testimony}.  Its content is demonstrably
derived from Hebrew Bible templates rather than from witness
testimony, oral tradition traced to an event, or documentary record.
The Gospels cannot therefore serve as evidence \textit{for} the
historicity of the crucifixion; they can serve only as evidence for
the fact that late first-century communities constructed
theologically meaningful narratives of a crucifixion using available
scriptural materials.

Whether any historical event lies beneath those constructions must be
established, if at all, from sources that are not themselves
scriptural derivations.  Paul might be such a source, and this paper
argues in §3 that the Pauline description of the crucifixion is
generic, soteriological, and non-biographical.  Josephus might be such
a source, and §6 argues that his attestation is weaker than is
usually claimed.  The argument is cumulative: templating neutralises
the Gospels as independent testimony; the remaining strata must then
carry the whole historicist burden.  They cannot carry that burden
without assistance from the Gospels they are being asked to
corroborate.  This is the sense in which Gospel templating is
important: it breaks the historicist cumulative case, not by
disproving the crucifixion, but by removing the strongest putatively-
independent evidence for it.

\paragraph{The criterion of embarrassment.}
The ``criterion of embarrassment'' is often invoked here: no one
would invent a crucified Messiah, so the crucifixion must be
historical.  This argument is addressed in detail in
Section~\ref{sec:myth:crucified}.  The short response is that
embarrassment does not on its own imply historicity---mythological
traditions routinely include transgressive or shocking elements that
are theologically generative rather than historically conceded---and
that the specific choice of crucifixion (as distinct from divine death
in general) is well-explained by the scriptural-curse logic of
Deuteronomy~21:23 as exploited by Paul in Galatians~3:13.  This does
not fully dissolve the argument, which retains some force; it does
show that ``no one would invent a crucified Messiah'' is not the
decisive move it has traditionally been taken to be.


%% ============================================================
\section{Non-Canonical Christian Sources}
\label{sec:noncanon}
%% ============================================================

The canonical Gospels are not the only early Christian texts.  A number
of non-canonical sources are relevant to the historicity question,
particularly because they preserve traditions that may be independent of
the canonical narrative.  The pattern that emerges from these sources
reinforces the findings of the preceding sections: the earliest
Christian texts construct their Jesus from scripture and revelation, not
from biographical memory.

%% ------------------------------------------------------------
\subsection{The Gospel of Thomas}
\label{sec:noncanon:thomas}
%% ------------------------------------------------------------

The Gospel of Thomas, preserved in a Coptic translation found at
Nag Hammadi in 1945, is a collection of 114 sayings attributed to
Jesus.  It contains no narrative framework: no birth, no ministry, no
miracles, no passion, no death, no resurrection.  It is a pure sayings
collection.

The dating of Thomas is contested.  Some scholars argue for a first-century
date for at least part of the collection, potentially independent of
the Synoptic Gospels \citep{Crossan1991,Patterson1993}.  Others place
it later and see dependence on the canonical Gospels
\citep{Goodacre2012}.

For the historicity question, the relevant observation is structural.
If Thomas preserves an early, independent tradition, that tradition
consists of wisdom sayings with no biographical frame---precisely the
kind of material that can be attributed to a personified Wisdom figure,
a divine revealer, or a philosophical school without requiring a
historical founder.  Many of the sayings in Thomas have parallels in
Jewish wisdom literature and Hellenistic philosophy.  The absence of
any passion or resurrection tradition in Thomas suggests either that the
community behind it did not know the narrative (consistent with the
mythicist model of progressive historicization) or that it deliberately
excluded it (requiring a theological explanation that itself undermines
the ``universal early tradition'' claim of historicists).

%% ------------------------------------------------------------
\subsection{The Didache}
\label{sec:noncanon:didache}
%% ------------------------------------------------------------

The Didache (``Teaching of the Twelve Apostles'') is an early Christian
manual of community practice, often dated to the late first or early
second century.  It contains instructions on baptism, fasting, prayer,
and the Eucharist, along with ethical teaching in the ``Two Ways''
tradition.

The Didache refers to ``the Lord'' as a source of teaching and
authority, but its references are ambiguous between a historical teacher
and a divine patron.  The ethical material in the ``Two Ways'' section
has Jewish parallels predating Christianity.  The eucharistic prayers
(chs.~9--10) give thanks over bread and wine but contain no reference
to the Last Supper, no institution narrative, and no connection to
a death or crucifixion.  The prayer gives thanks for ``the life and
knowledge which you made known to us through Jesus your servant''
(\grk{pais})---language that can refer to a revelatory intermediary as
easily as to a historical teacher.

The Didache, like Thomas, presents a Jesus tradition that is functional
and liturgical rather than biographical.  It presupposes a Lord and a
teaching, not a life story.

%% ------------------------------------------------------------
\subsection{1~Clement and the Epistle of Barnabas}
\label{sec:noncanon:clement}
%% ------------------------------------------------------------

1~Clement, a letter from the Roman church to the Corinthian church,
is typically dated to the mid-90s~CE.  It is therefore roughly
contemporary with the later canonical Gospels.  Its treatment of Jesus
is striking.

Clement discusses Christ extensively but draws his knowledge of Christ
almost entirely from Hebrew Bible citations and early Christian creedal
formulae.  When Clement comes to describe Jesus's suffering (ch.~16),
he does not paraphrase or cite any Gospel account; he reproduces
Isaiah~53 and Psalm~22 directly.  As \citet{Doherty2009} observes,
Clement appears to derive his knowledge of the passion from scripture,
not from any narrative tradition about historical events.

Clement's Jesus speaks---but the words attributed to him are quotations
from the Psalms and other Hebrew Bible texts, presented as words of
Christ spoken through the prophets.  This is consistent with an
understanding of Christ as a pre-existent divine figure who spoke
through scripture, not as a recently-living teacher whose words were
transmitted by memory.

Clement's silence about Gospel traditions is comprehensive.  He
discusses humility, obedience, and proper order without citing the
Sermon on the Mount.  He addresses disputes in the Corinthian community
without citing any dominical sayings.  He refers to the apostles as
appointed by Christ but provides no narrative of their calling.  He
mentions the resurrection but offers no empty-tomb narrative.  He
is writing to Corinth---Paul's own community---and yet his knowledge of
Christ is mediated entirely by scripture and creed, not by biographical
tradition.

The Epistle of Barnabas (typically dated to the early second century)
is even more revealing.  Barnabas constructs an elaborate Christology
entirely from Hebrew Bible typology: the scapegoat ritual
(Leviticus~16), the red heifer (Numbers~19), the Suffering Servant
(Isaiah~53), and a range of other scriptural types.  Jesus in Barnabas
is a figure constructed from exegesis.  There is no reference to
earthly events, named places, or historical contemporaries.  If one did
not know the Gospels existed, reading Barnabas would give no reason to
think a historical person was involved in the origins of the tradition
\citep{Doherty2009}.

%% ------------------------------------------------------------
\subsection{The Ascension of Isaiah}
\label{sec:noncanon:ascension}
%% ------------------------------------------------------------

The Ascension of Isaiah, a composite text whose Christian section
(chs.~6--11) may date to the late first or early second century,
describes a vision of ``the Beloved'' descending through the heavens and
being crucified by demonic powers who do not recognize him.
\citet{Doherty2009} and \citet{Carrier2014} treat this as evidence for
an early tradition in which the crucifixion was cosmic rather than
historical; critics \citep{McGrath2014} argue the text includes a
descent to earth and reflects Docetism rather than mythicism.  The
textual situation is genuinely complex, with multiple versions and
significant variation.

The Ascension of Isaiah is not load-bearing for the mythicist case.  It
provides a window into the diversity of early Christian cosmological
thought---a diversity in which a celestial, non-historical Christ was at
least conceivable---but the core mythicist argument rests on the Pauline
evidence, the silence of non-Christian sources, and the compositional
character of the Gospels, not on this contested text.


%% ============================================================
\section{Non-Christian Attestation}
\label{sec:external}
%% ============================================================

The non-Christian sources for Jesus are routinely cited as external
corroboration of the Gospel tradition.  There are remarkably few of
them, and each, on examination, is less probative than the consensus
maintains.

%% ------------------------------------------------------------
\subsection{Josephus: The \textit{Testimonium Flavianum}}
\label{sec:external:tf}
%% ------------------------------------------------------------

The most frequently cited non-Christian reference to Jesus is a passage
in Flavius Josephus's \textit{Antiquities of the Jews} (18.3.3),
written \circa 93--94~CE.  In its received form, the passage reads:

\begin{quote}
About this time there lived Jesus, a wise man, \textit{if indeed one
ought to call him a man}.  For he was one who performed surprising
deeds and was a teacher of such people as accept the truth gladly.  He
won over many Jews and many of the Greeks.  \textit{He was the
Messiah.}  And when, upon the accusation of the principal men among us,
Pilate had condemned him to a cross, those who had first come to love
him did not cease.  \textit{He appeared to them spending a third day
restored to life, for the prophets of God had foretold these things and
a thousand other marvels about him.}  And the tribe of the Christians,
so called after him, has still to this day not disappeared.
\end{quote}

\noindent
The italicized phrases are almost universally recognized as Christian
interpolations \citep{Meier1991,Ehrman2012}.  Josephus, a non-Christian
Jew writing for a Roman audience, would not have called Jesus the
Messiah, would not have affirmed the resurrection, and would not have
described him as more than human.

The scholarly debate concerns whether, after removing the interpolated
phrases, a genuine Josephan core remains.  The ``partial authenticity''
position---that Josephus wrote a shorter, neutral reference to Jesus
which Christian scribes subsequently embellished---is the current
majority view \citep{Meier1991}.  Under this reconstruction, the
original passage might have read something like: ``About this time there
lived Jesus, a wise man.  He was a teacher\ldots\ Pilate condemned him
to a cross\ldots\ the tribe of the Christians has not disappeared.''

This reconstruction, however, is not an observation; it is a
\textit{hypothesis} about what Josephus might have written.  No
manuscript preserves the uninterpolated text.  The ``core'' is reverse-
engineered by removing what seems Christian and retaining what seems
neutral.  The result is a passage whose existence is inferred, not
attested.

Several considerations weigh against even the partial authenticity
view:

\begin{itemize}[leftmargin=2em]
\item \textbf{Narrative flow.}  Remove the entire \textit{Testimonium}
and the surrounding text reads more naturally.  The passage immediately
before discusses a disturbance involving Pilate, and the passage
immediately after begins ``About the same time also another sad
calamity put the Jews into disorder\ldots''  The word ``another''
(\grk{heteron}) connects to the prior Pilate episode, not to the
\textit{Testimonium}, suggesting the \textit{Testimonium} was inserted
into an existing sequence \citep{Olson1999}.

\item \textbf{Origen's silence.}  Origen, writing in the mid-third
century, explicitly states that Josephus did not believe Jesus was
the Messiah (Contra Celsum~1.47).  This is consistent with either
an original neutral reference or no reference at all.  But Origen
also appears to cite Josephus as connecting the fall of Jerusalem
with the death of James---a connection not found in our current text
of Josephus.  This suggests the text of Josephus available to Origen
differed from what we now possess, raising the possibility of
significant editorial intervention between Origen's time and the
earliest surviving manuscripts.

\item \textbf{Eusebius.}  The first author to quote the
\textit{Testimonium} in its full received form is Eusebius of
Caesarea (\circa 324~CE).  \citet{Olson1999} has argued that Eusebius
himself may be the interpolator, noting that Eusebius explicitly
endorsed the use of pious fraud in the service of the faith
(\textit{Praeparatio Evangelica}~12.31).

\item \textbf{The argument from brevity.}  If Josephus wrote a neutral
reference to Jesus, it is strikingly brief for a writer who devotes
extensive space to other Jewish religious figures, messianic claimants,
and prophetic movements of the period.  Josephus's account of John the
Baptist (\textit{Antiquities}~18.5.2) is substantially longer and more
detailed.  A brief, passing reference---if genuine---would suggest that
Jesus was, in Josephus's assessment, a minor figure, which contradicts
the Gospel portrayal but is consistent with a late editorial insertion
rather than an organic part of the narrative.
\end{itemize}

The honest assessment is that the \textit{Testimonium Flavianum} cannot
be used as evidence for the historical Jesus.  In its received form it
is partially or wholly interpolated.  The hypothesized neutral core is
a modern reconstruction, not a preserved text.  And even if a genuine
Josephan reference existed, it would tell us only that Josephus, writing
sixty years after the supposed events, was aware of Christians and their
claims about their founder---which is unsurprising and provides no
independent corroboration of those claims.

%% ------------------------------------------------------------
\subsection{Josephus: The James Passage}
\label{sec:external:james}
%% ------------------------------------------------------------

The second Josephan reference (\textit{Antiquities}~20.9.1) is
generally considered more secure.  Describing the illegal execution of
certain men by the high priest Ananus in 62~CE, Josephus writes that
Ananus ``convened the judges of the Sanhedrin and brought before them
the brother of Jesus who is called Christ [\grk{ton adelphon I\=esou
tou legomenou Christou}], James by name, and certain others.''

This passage is less obviously interpolated: the reference to Jesus
serves to identify which James is meant, and the phrase ``who is called
Christ'' (\grk{tou legomenou Christou}) is neutral or even skeptical,
which argues against Christian insertion.

However, several problems remain:

\begin{itemize}[leftmargin=2em]
\item \textbf{Dependence on the \textit{Testimonium}.}  The phrase
``Jesus who is called Christ'' uses the identifying tag as though
Jesus has already been introduced.  If the \textit{Testimonium} in
Book~18 is wholly interpolated, then Book~20 refers back to a person
Josephus never mentioned---which suggests that the Book~20 reference
may also be an insertion, or at minimum that the identifying clause
\grk{tou legomenou Christou} was added by the same hand that
inserted the \textit{Testimonium} \citep{Carrier2014}.

\item \textbf{Possible gloss.}  The phrase ``the brother of Jesus who
is called Christ'' could be a marginal gloss that migrated into the
text.  A Christian scribe copying the passage, recognizing that this
James was the James known in Christian tradition, could have added
the identifying tag in the margin, which a subsequent copyist
incorporated.  This is the same mechanism proposed for the
\grk{ton adelphon tou kyriou} phrase in Galatians~1:19
(Section~\ref{sec:paul:brother}).

\item \textbf{Origen's discrepancy.}  As noted above, Origen cites
Josephus as connecting the fall of Jerusalem with the death of
James---but the passage in our \textit{Antiquities} makes no such
connection.  This suggests either that Origen was citing a different
passage (now lost), that the text has been altered between Origen's
time and our manuscripts, or that Origen was misremembering.  Any of
these possibilities complicates the assumption that the text we
possess faithfully represents what Josephus wrote.
\end{itemize}

The honest assessment of the James passage is that it is genuinely
ambiguous.  The majority of textual scholars accept the passage as
substantially authentic, and the arguments for interpolation are
\textit{possible} rather than \textit{probable}.  The brief, neutral,
non-Christianizing character of the reference (``who is called Christ''
is skeptical or at most neutral in tone) argues against wholesale
Christian insertion.

If the passage is authentic, it establishes that Josephus, writing in
the 90s~CE, identified a man named James by reference to a figure called
Christ---and did so in a way that suggests he expected his readers to
know who this Christ was, presumably from his earlier discussion
(whether that discussion was the \textit{Testimonium} or something
else).  This is evidence---not conclusive, but genuine---that by the
late first century, the connection between a historical James and the
Christ figure was known outside exclusively Christian circles.  The
historicist reads this as confirming that Jesus existed and had a
brother named James.  The mythicist reads it as confirming that the
historicized narrative was in circulation by the 90s, which is
consistent with post-Markan development.  The passage cannot decisively
distinguish between these readings, though it mildly favors historicity
in isolation.  Its force, however, is limited by the dependence problem
noted above: if the \textit{Testimonium} is wholly interpolated, the
Book~20 reference may be as well.

%% ------------------------------------------------------------
\subsection{Tacitus, \textit{Annals}~15.44}
\label{sec:external:tacitus}
%% ------------------------------------------------------------

Writing \circa 116~CE, the Roman historian Tacitus describes Nero's
persecution of Christians following the great fire of Rome in 64~CE.
He notes that the name ``Christian'' derives from ``Christus,'' who
``suffered the extreme penalty during the reign of Tiberius at the hands
of one of our procurators, Pontius Pilatus.''

This passage is frequently cited as independent Roman confirmation of
Jesus's crucifixion under Pilate.  Several considerations diminish its
evidential value:

\begin{itemize}[leftmargin=2em]
\item \textbf{Date.}  Tacitus is writing more than eighty years after
the supposed crucifixion.  His information is not contemporary.

\item \textbf{Source.}  Tacitus does not cite a Roman archival source.
The passage contains the title ``procurator'' (\textit{procurator})
for Pilate, whereas Pilate's actual title was ``prefect''
(\textit{praefectus}), as confirmed by the Pilate inscription
discovered at Caesarea Maritima in 1961 \citep{Vardaman1962}.  The
use of the wrong title suggests Tacitus is not working from official
records, which would have used the correct designation.  He is more
likely reporting what was common knowledge about Christians in his
time---\ie, what Christians themselves said about their origins.

\item \textbf{Christian source contamination.}  By 116~CE, Christian
communities were well established in Rome.  Tacitus's information
about the origin of the name ``Christian'' and the circumstances of
the founder's death almost certainly derives, directly or indirectly,
from Christian tradition.  The passage tells us what Christians
believed, not what independently verified Roman records contained.

\item \textbf{Textual questions.}  The \textit{Annals} survive in a
single manuscript tradition.  The earliest manuscript containing
Book~15 dates to the eleventh century.  While wholesale fabrication
of the Tacitus passage is considered unlikely by most scholars, the
possibility of minor interpolation or editorial modification over a
millennium of transmission cannot be excluded.
\end{itemize}

The Tacitus passage confirms that by the early second century a Roman
historian---one known for care and hostility toward his subjects---was
aware that Christians claimed their founder had been executed under
Pilate.  It does not provide independent confirmation that this
execution occurred, because Tacitus's source is almost certainly the
Christian tradition itself.  However, the passage is not valueless: it
establishes that the historicized narrative (execution under Pilate) was
in circulation early enough to reach a Roman historian writing
\circa 116~CE, and that Tacitus saw no reason to question or qualify it.
Under the mythicist model, this means the historicization process was
substantially complete by the early second century---which is consistent
with a timeline in which Mark (\circa 70~CE) introduced the Palestinian
narrative framework, but which places constraints on how late the
historicization can be dated.

%% ------------------------------------------------------------
\subsection{Pliny, Suetonius, and Other Minor References}
\label{sec:external:minor}
%% ------------------------------------------------------------

Pliny the Younger, writing to the emperor Trajan \circa 112~CE, describes
Christians in Bithynia who sing hymns to ``Christus as to a god''
(\textit{Epistulae}~10.96).  This tells us that Christians worshipped
Christ as a deity in the early second century.  It tells us nothing about
whether the figure they worshipped was historical.

Suetonius, in his \textit{Life of Claudius} (25.4), mentions that
Claudius expelled Jews from Rome because of disturbances ``at the
instigation of Chrestus'' (\textit{impulsore Chresto}).  The spelling
``Chrestus'' may refer to Christ or may refer to an otherwise unknown
individual named Chrestus---a common slave name.  Even on the Christian
reading, the passage tells us only that disputes about Christ were
causing disturbances in the Roman Jewish community \circa 49~CE, not
that Christ was a historical person.

The Talmudic references to Jesus (primarily in \textit{b.~Sanhedrin}
43a and 107b) date from the third to fifth centuries~CE and are
transparently shaped by polemical engagement with Christian claims
rather than by independent historical tradition.  \citet{Ehrman2012}
himself acknowledges that this material is too late to be of significant
historical value.

Mara bar Serapion, a Syrian Stoic writing from prison sometime after
73~CE (the date is disputed; some place the letter much later), refers
to the Jews executing ``their wise king.''  The identification with
Jesus is not certain---the reference could be to another figure---and
even if it refers to Jesus, the information is vague and likely derived
from Christian tradition rather than independent knowledge.

In sum, the non-Christian attestation for Jesus consists of: one
partially or wholly interpolated passage in Josephus; one brief
identifying tag in Josephus whose authenticity is debated but which the
majority of scholars accept; one reference in Tacitus almost certainly
derived from Christian sources but confirming early circulation of the
historicized narrative; and a handful of minor references that tell us
about the existence of Christians and their beliefs.

This inventory is weaker than the consensus typically claims, but it is
not empty.  The non-Christian sources collectively establish that by the
late first and early second centuries, the claim that Christianity had
a historical founder executed under Pilate was in wide circulation and
was accepted even by hostile or neutral observers.  Under the historicist
model, this confirms a historical core.  Under the mythicist model, it
confirms that the historicization process was rapid and effective.  The
non-Christian evidence alone cannot decisively distinguish between these
readings, but it does constrain the mythicist timeline and establishes
that the historicized tradition was not a late development.

%% ------------------------------------------------------------
\subsection{Rabbinic References and \textit{Toledoth Yeshu}}
\label{sec:external:rabbinic}
%% ------------------------------------------------------------

The Talmud and related rabbinic literature contain several passages
that mention a figure variously named \grk{Yeshu}, \grk{Yeshu
ha-Notzri} (\enquote{Yeshu the Nazarene}), or \grk{Yeshu ben Pandera}.
These are sometimes invoked as non-Christian attestation of the
historical Jesus.  The textual and interpretive state of this material
is rigorously surveyed by \citet{SchaferP2007} and by
\citet{MeersonSchafer2014}; the following summarises the principal
cases and their evidential weight.

\textbf{b.Sanhedrin 43a.}  \enquote{On the eve of Passover, Yeshu
was hanged.  For forty days beforehand a herald went out in front
of him, crying: `He is going out to be stoned for practicing magic
and leading Israel astray\ldots'}  This is the most specific
talmudic reference and has sometimes been taken as independent
rabbinic memory.  Schäfer argues it is a late (post-third-century)
Babylonian composition that inverts the Christian passion narrative
for polemical effect: stoning rather than crucifixion, the formal
due-process characterisation serving as counter-claim against
Christian allegations of an unjust execution.  The passage assumes
knowledge of the Christian narrative rather than providing
independent corroboration of it.

\textbf{b.Sanhedrin 107b and b.Sotah 47a.}  Yeshu is described as a
disciple of Rabbi Joshua ben Perachiah who was excommunicated for
apostasy and idolatry.  The chronological anachronism is large
(Joshua ben Perachiah lived in the early first century BCE, not
the first century CE); the passage is generally treated as
polemical construction, not biographical memory.

\textbf{b.Shabbat 104b and b.Sanhedrin 67a.}  Yeshu ben Pandera /
Yeshu ben Stada: a magician who brought spells from Egypt.  The
\enquote{ben Pandera} tradition reappears in Origen's
\textit{Contra Celsum} (c.~248~CE) where Celsus's Jewish
informant attributes a similar origin story to Jesus.  This is the
earliest documented rabbinic-tradition material reaching a
Christian writer.

\textbf{\textit{Toledoth Yeshu}.}  The \textit{Toledoth Yeshu}
compilations (existing in many manuscript recensions; the earliest
evidence for oral traditions is patristic, with Agobard of Lyon
c.~830 describing the general form) present an extended anti-
Christian counter-narrative: Yeshu as illegitimate son of Miriam
and a Roman soldier named Pandera or Pantera, trained as a
magician, performing miracles by magical means, executed, and
buried in a disgraceful grave.  \citet{MeersonSchafer2014} provide
the critical edition and argue that these compilations
consolidate oral polemical traditions of varying ages, but that
the tradition as a whole originates as \textit{polemical response
to the Christian gospel narrative}, not as independent memory of a
historical figure.

\paragraph{Evidential assessment.}
None of the rabbinic passages provides early, independent, non-
polemical attestation of a first-century Galilean preacher.  All of
them either post-date the Christian narrative by centuries or
are most plausibly understood as polemical engagement with it.
Their evidential value for historicity is therefore similar to
Tacitus's: they confirm that Christian claims about a historical
founder were circulating early enough to attract Jewish polemical
response by the late second or early third century, and they
inherit the historicizing frame of the Christian tradition without
providing independent witness to its content.  They do not
establish historicity, and they do not disconfirm it either; like
the Josephan and Tacitean material, they constrain both hypotheses
to accommodate the early circulation of a historicized narrative
without favouring either side of the existence question.

A note on \textit{Toledoth Yeshu}'s significance for the paper's
argument.  The \textit{Toledoth Yeshu} material has been read by
some mythicist-sympathetic writers as evidence that the earliest
Jewish engagement with the Christian claim treated it as a
fabrication rather than as a memory of a known figure.  This
reading is over-stated: polemical counter-narratives routinely
invert the claims they oppose without thereby taking a position on
the existence of the underlying subject, and the ben Pandera
tradition assumes enough of the Christian figure to count as
engagement rather than denial.  What the tradition does establish
is that the rabbinic community's earliest substantial response to
Christianity was polemical reshaping of the Gospel narrative, not
independent testimony to events.

%% ------------------------------------------------------------
\subsection{The {Mandaean} Tradition}
\label{sec:external:mandaean}
%% ------------------------------------------------------------

The Mandaeans are a surviving religious community of Southern
Mesopotamian origin with demonstrably pre-modern roots and a
substantial textual tradition in Eastern Aramaic (Mandaic).  They are
centred on John the Baptist (Mandaic \grk{Yahya} or \grk{Yuhana})
rather than on Jesus, and their canonical literature contains
independent references to Jesus that are almost entirely ignored in
the historical-Jesus scholarly literature.  \citet{Rudolph1978},
\citet{Lupieri2002}, and \citet{BuckleyMandaeans2002} give the
standard recent surveys.

\paragraph{The community and its sources.}
The Mandaean canonical texts include:
\begin{itemize}[leftmargin=2em]
\item the \textbf{Ginza Rabba} (\enquote{Great Treasure}), a
compendium of mythological, theological, liturgical, and
historical material;
\item the \textbf{Drasha d-Yahya} (\enquote{Teaching of John}, or
\enquote{Book of John});
\item the \textbf{Qolasta} (the canonical prayerbook);
\item the \textbf{Haran Gawaita}, a short historical-mythological
text that narrates the community's migration from Palestine/Judea
to Mesopotamia under persecution.
\end{itemize}
Dating is the central problem.  The manuscripts we possess are
late medieval; the texts themselves are generally agreed to contain
layered traditions of varying ages, with some scholars
(\citealt{Lidzbarski1925}, \citealt{BuckleyMandaeans2002}) arguing
for first- or second-century origins for portions of the material,
and others (\citealt{Lupieri2002}) for a later Mesopotamian
composition with earlier traditions embedded.  The \textit{Haran
Gawaita}'s self-report of a first-century CE westward origin is
contested but is not in itself implausible.

\paragraph{Mandaean references to Jesus.}
Mandaean texts refer to Jesus as \grk{Yi\v{s}u Mshiha} and treat him
as a \textit{false} messiah---a corrupter of John the Baptist's
teaching, associated with the evil spirit \grk{R\^uha} and at times
with the planet Mercury in Mandaean astrological symbolism.  The
principal passages are scattered through the Ginza Rabba
(especially Book~1, Right Ginza) and the Book of John.  The Jesus
figure is given narrative detail: his baptism by John (accepted by
the Mandaean tradition but with reversed valuation, Jesus as the
one baptized who then corrupts the lineage), his claim to
messiahship, and his execution.  He is treated as a specific
historical figure whose existence is not denied but whose
significance is inverted.

\paragraph{Evidential assessment.}
The Mandaean tradition is potentially the most interesting
unexploited witness for the historicity question, but the
evidential gain is limited by what the tradition actually preserves
and by the dating problem.  Three considerations apply.

\textit{First, the Mandaean Jesus is not a biographically detailed
figure.}  The references assume a Jesus whose existence is the
framework for polemical reversal of Christian claims, but they do
not preserve biographical content that is independent of the
Christian narrative itself.  Like the rabbinic material in
§\ref{sec:external:rabbinic}, the Mandaean Jesus appears to be
known through polemical engagement with the Christian tradition
rather than through independent memory of a first-century figure.

\textit{Second, the dating problem is severe.}  The Mandaean texts
in their current form post-date Christianity by centuries; even the
most generous reconstructions of the earlier oral traditions
cannot rule out that the Mandaean Jesus-material originated as
response to Christian claims circulating in Mesopotamia by the
third or fourth century.

\textit{Third, and crucially, the Mandaean community's existence
is itself a datum.}  If (as the \textit{Haran Gawaita} claims and
as some scholarship supports) the Mandaeans derive from a
first-century Palestinian John-the-Baptist-aligned sectarian
community that moved east under Christian and Roman pressure,
then the Mandaean tradition represents a living continuation of
the Second Temple Jewish apocalyptic-mystical milieu in which
earliest Christianity formed.  Whether the Mandaean Jesus is
independent memory or later polemical construction, the Mandaean
existence is evidence for the diversity of the religious milieu
from which Christianity differentiated---and that diversity
constrains any theory of the earliest Christian community's
emergence.

\paragraph{Honest status.}
The Mandaean material \textit{does not provide} independent
first-century attestation of a historical Jesus; the tradition is
too late and too dependent on the Christian narrative to carry that
weight.  It \textit{does not refute} historicity either.  What it
offers is a live pre-Islamic (and perhaps pre-Christian-in-origin)
religious tradition that knew Jesus as a contested polemical
figure rather than as the centre of its own theology.  A more
thorough engagement with the Mandaean material---specifically, an
analysis of whether any Mandaean Jesus-references can be textually
dated to pre-Constantinian Mesopotamian material---would constitute
genuinely novel evidence for either side of the debate.  The field
has not done this work; we note the opportunity without claiming to
have done it here.

%% ------------------------------------------------------------
\subsection{Archaeological Substrate: {Nazareth}, Sepphoris, and the Yehohanan Ossuary}
\label{sec:external:archaeology}
%% ------------------------------------------------------------

Archaeology occasionally settles historicity questions decisively:
the 1961 Pilate inscription at Caesarea Maritima \citep{Vardaman1962}
moved Pilate from literary-only attestation to physical attestation.
For Jesus, no comparable artefact exists.  Four archaeological
considerations nonetheless bear on the question.

\paragraph{The habitation of first-century {Nazareth}.}
\citet{Salm2008} argued a skeptical case that Nazareth was
uninhabited in the first century CE, making the village setting of
the Gospel nativity and youth narratives a literary fiction.  His
argument rested on the alleged absence of datable first-century
remains in early Israeli excavations.  \citet{Reed2000},
\citet{Dark2020}, and specialist archaeologists (notably James
Strange and Stephen Pfann) have responded in detail, arguing that
first-century CE habitation is attested by pottery, rock-cut tombs
(e.g., Kokhim tombs of Jewish type), agricultural installations
(wine and olive presses), and some structural remains recently
re-examined by the Nazareth Archaeological Project.  The mainstream
archaeological consensus now is that Nazareth was a small Jewish
village through the first century, though with a very modest
footprint.  Salm's case has been effectively answered in the
specialist literature.

The implication for historicity is modest.  A village Nazareth's
existence does not confirm Jesus.  A non-existent Nazareth would
have been a substantial argument against the Gospel narratives, and
that argument is not available.  The Gospel-era Nazareth was a
small, obscure village---which is neither surprising nor
informative on the principal question.

\paragraph{The silence about Sepphoris.}
Sepphoris was a major Hellenistic-Jewish urban centre approximately
four miles (six kilometres) from Nazareth, serving as the
administrative capital of Galilee under Herod Antipas.  By any
plausible reconstruction of a Nazareth-based Galilean ministry, a
public-teaching figure would have had to engage with Sepphoris or
deliberately to avoid it.  \textit{Sepphoris is conspicuously absent
from every Gospel text.}  This is a minor datum but worth
registering: the Gospel narratives construct their Galilean setting
around small villages (Nazareth, Capernaum, Cana, Bethsaida) and
never mention the regional administrative capital.  If the Gospel
writers were reconstructing a historical ministry from memory, this
absence is puzzling; if they were constructing a narrative from
scriptural Galilean imagery (fishermen, villages, boats, grain
fields) without specific historical constraint, the absence is
expected.  Sepphoris is not a decisive datum, but it is a
consilience pointing in the paper's direction.

\paragraph{The {Yehohanan} ossuary and the archaeology of
crucifixion.}
\citet{TzaferisYehohanan1970} and \citet{ZiasSekeles1985} published
and reappraised the bone remains of Yehohanan ben Hagkol, a first-
century CE crucifixion victim whose heel bone was discovered in
1968 with a Roman iron nail still in place.  This is the \textit{only}
surviving archaeological evidence of Roman crucifixion in Judea.
Its relevance to historicity is this: crucifixion was a real
Roman practice in first-century Palestine---we have physical
evidence---but the practice left almost no archaeological trace in
general, because bodies were typically disposed of in mass or open
pits and because nails of this kind were routinely reused.  The
absence of an archaeological trace for a specific crucified Jesus
is thus of negligible evidential weight either way.

\paragraph{The James ossuary and forgery-determination.}
In 2002 an ossuary inscribed with the Aramaic phrase
\enquote{James, son of Joseph, brother of Jesus}
(\grk{Ya'akov bar Yosef achui d'Yeshua}) became briefly a
candidate for the first direct archaeological attestation of
Jesus.  The subsequent investigation by the Israel Antiquities
Authority and ensuing trial (Oded Golan, 2005--2012) concluded that
at least the \enquote{brother of Jesus} portion of the inscription
was a modern forgery added to an otherwise genuine first-century
ossuary.  The case is legally concluded but the physical and
epigraphic analysis remains contested among specialists.  The
episode is instructive about the field's sensitivity to any
prospect of physical attestation---and, negatively, about how
thoroughly the true archaeological record has been searched.  No
genuine artefact mentioning Jesus from the first century has been
recovered, and the field's willingness to treat a contested
inscription as potentially definitive evidence signals the evidentiary
pressure the historicity question is under.

\paragraph{Summary.}
Archaeology provides no direct evidence for Jesus.  The
surrounding contextual archaeology (first-century Nazareth's
existence, Pilate's confirmed prefecture, the physical reality of
crucifixion in Judea) is consistent with the Gospel setting but
does not confirm a specific figure.  The absence of direct
physical attestation is not decisive in itself---many first-
century Galilean peasants are archaeologically invisible---but it
places Jesus on the Socrates/Pythagoras side of the evidential
profile rather than the Pilate/Augustus side.


%% ============================================================
\section{The Argument from Silence}
\label{sec:silence}
%% ============================================================

Section~\ref{sec:method:silence} established the formal principle that
absence of expected evidence is evidence of absence, with strength
proportional to how confidently the evidence would be expected under the
hypothesis.  This section applies that principle to specific documentary
traditions that should contain references to Jesus if the historicist
hypothesis is correct.

%% ------------------------------------------------------------
\subsection{Philo of Alexandria}
\label{sec:silence:philo}
%% ------------------------------------------------------------

Philo Judaeus (\circa 20~BCE--\circa 50~CE) was a Hellenistic Jewish
philosopher who wrote extensively on Jewish theology, law, philosophy,
and politics.  His surviving corpus is enormous---over forty treatises.
He was an exact contemporary of the supposed Jesus and lived in
Alexandria, a major center of Jewish intellectual life closely connected
to Jerusalem.

Philo's writings are directly relevant to the historicity question on
multiple fronts:

\begin{itemize}[leftmargin=2em]
\item He discusses Pontius Pilate at length, describing his
administration and its conflicts with the Jewish population
(\textit{Legatio ad Gaium}~299--305).

\item He writes extensively about Jewish messianic expectation, about
the interpretation of Hebrew Bible prophecy, about Wisdom and Logos
theology---precisely the intellectual milieu from which Christianity
emerged.

\item He describes Jewish sectarian movements, including the Essenes
and the Therapeutae.

\item He lived during the period when Jesus supposedly taught,
attracted followers, was executed, and his movement began spreading
across the Mediterranean.
\end{itemize}

Philo never mentions Jesus, Christians, or anything recognizable as the
Christian movement.

The historicist response is that Jesus was too obscure to attract
Philo's notice.  This defense has more force than mythicists typically
acknowledge.  Philo lived in Alexandria, not Jerusalem.  He was a
philosopher, not a chronicler of Galilean preachers.  His writings
are thematically focused on scriptural interpretation, Jewish law, and
relations with Roman authority---topics that would not necessarily lead
him to discuss an itinerant teacher in a distant province.  If the
historical Jesus was a minor figure---a local teacher with a small
following, executed as one of many would-be messiahs---Philo's silence
is unremarkable.

However, this defense has limits.  If a significant religious movement
emerged from within Palestinian Judaism during Philo's lifetime---a
movement making extraordinary claims about a recently-executed
Messiah---and if this movement was sufficiently organized to produce
communities across the Mediterranean within two decades (as Paul's
letters attest), then its complete absence from the most prolific Jewish
writer of the period is at least noteworthy.  The strength of this
silence as evidence depends on how visible the early Christian movement
was: if it was a tiny, obscure sect (as historicists who invoke the
``obscurity'' defense must maintain), Philo's silence is unsurprising;
but this smallness claim is in tension with the Gospel narratives
describing large crowds and public confrontation.

Under the historicist hypothesis, $P(\text{Philo silent} \mid H)$ is
difficult to estimate---it depends on how prominent Jesus and his
movement actually were, which is precisely what is in dispute.  The
silence of Philo is suggestive but weaker than a maximalist mythicist
reading would claim.

%% ------------------------------------------------------------
\subsection{Justus of Tiberias}
\label{sec:silence:justus}
%% ------------------------------------------------------------

Justus of Tiberias was a Jewish historian and contemporary of Josephus
who wrote a history of the Jewish war and, importantly, a chronicle of
events in Galilee---the region where Jesus supposedly conducted his
ministry.  His works are lost, but the ninth-century patriarch Photius
read them and recorded his reaction in his \textit{Bibliotheca}
(Codex~33):

\begin{quote}
[Justus] makes not the least mention of the appearance of Christ, or
what things happened to him, or of the wonderful works that he did.
\end{quote}

\noindent
Photius, a Christian reader, found this silence remarkable.  And well
he might: a historian of Galilee, writing about the very region and
period of Jesus's supposed activity, who makes no mention of Jesus or
his followers.

The historicist defense is that Justus's work is lost and we cannot
know its full contents from Photius's brief notice.  This is true but
irrelevant: Photius specifically comments on the absence of any mention
of Christ, indicating that he searched for one and did not find it.
Given that Photius was a Christian patriarch with every motivation to
find such a reference, his explicit notice of its absence carries
significant weight.

%% ------------------------------------------------------------
\subsection{The Mishnaic Tradition}
\label{sec:silence:mishnah}
%% ------------------------------------------------------------

The Mishnah, compiled \circa 200~CE from earlier oral traditions,
contains extensive discussions of Sanhedrin procedure, heresy trials,
capital punishment, the authority of the high priest, and the
relationship between Jewish and Roman legal jurisdiction---precisely
the institutional context in which Jesus supposedly appeared, was tried,
and was condemned.

The Mishnah contains no recognizable reference to the trial or
execution of Jesus.  However, the evidential weight of this silence must
be carefully calibrated.  The Mishnah is a legal code, not a chronicle
of historical events.  It records legal principles and procedural norms,
not specific cases.  It contains no reference to the trial of Hillel,
or to any specific execution outside its legal hypotheticals.  The
absence of Jesus's trial from the Mishnah is therefore less surprising
than it might initially appear---the Mishnah omits most specific
historical events.

That said, if the Gospel trial narrative is historical, it describes
proceedings that egregiously violated Mishnaic procedural norms (a night
trial, a capital case decided in a single day, a trial during a
festival).  The Mishnah's framers might have had reason to address such
a case precisely because of its procedural irregularity.  The silence
is mildly suggestive but not strong evidence on its own.

The later Talmudic references (discussed in
Section~\ref{sec:external:minor}) are too late and too obviously
polemical to represent independent tradition.  They respond to
Christian claims rather than preserving pre-Christian memory.

%% ------------------------------------------------------------
\subsection{Roman Administrative Records}
\label{sec:silence:roman}
%% ------------------------------------------------------------

The Roman Empire maintained administrative records of provincial
governance, including judicial proceedings.  While the vast majority of
such records have been lost, the complete absence of any reference to a
trial or execution of Jesus in any surviving Roman document, inscription,
or administrative text is notable---particularly because we do possess
Roman references to other events in Judean governance from the same
period.

The Pilate inscription from Caesarea Maritima confirms Pilate's
historicity and his title (\textit{praefectus}).  We have coins minted
under Pilate's authority.  We have Philo's and Josephus's independent
accounts of Pilate's administration.  The Roman apparatus in Judea left
traces.  The trial and execution of Jesus did not---unless one counts
the Tacitus passage, which, as argued above, derives from Christian
tradition rather than Roman records.

The force of this silence depends on how many records we would expect
to survive.  Individual provincial trials were unlikely to be preserved
in the historical record on their own.  But the Gospel accounts do not
describe an ordinary trial.  They describe a case involving the Roman
prefect personally, a Sanhedrin convened for the purpose, crowds
demanding execution, extraordinary portents (earthquake, darkness,
temple veil rent), and an executed man's body going missing.  If even a
fraction of this occurred, the event was extraordinary, and extraordinary
events leave traces.

%% ------------------------------------------------------------
\subsection{Cumulative Force of Absent Evidence}
\label{sec:silence:cumulative}
%% ------------------------------------------------------------

No single silence is decisive.  The force of the argument lies in its
cumulative structure.

The following documentary channels would plausibly contain a reference
to Jesus if the historicist hypothesis is true.  They are listed in
approximate order of expected sensitivity---\ie, the probability that
each would mention Jesus given historicity---from strongest to weakest:

\begin{enumerate}[leftmargin=2em]
\item \textbf{Justus of Tiberias} (late first century): a historian
of Galilee---Jesus's supposed home region---who made no mention of
Christ, as the ninth-century patriarch Photius explicitly noted.
This is among the strongest individual silences, though it depends on
a single late witness to a lost work.

\item \textbf{Philo of Alexandria} (\circa 20~BCE--50~CE): the most
prolific Jewish writer of the period, who discussed Pilate, messianic
expectation, and Jewish sectarian movements.  No mention.  The force
of this silence is moderated by Philo's location (Alexandria, not
Jerusalem) and his philosophical rather than chronicling orientation
(Section~\ref{sec:silence:philo}).

\item \textbf{Roman administrative records from Judea}: no surviving
reference to a trial or execution of Jesus, despite the survival of
other traces of Pilate's administration.  The force is moderated by
the general loss of provincial records from this period.

\item \textbf{Seneca the Younger} (\circa 4~BCE--65~CE): Stoic
philosopher and tutor to Nero, who wrote on Jewish topics.  No
mention.

\item \textbf{The Mishnaic legal tradition}: compiled \circa 200~CE
from earlier oral traditions.  No reference to a trial of Jesus.  The
force is weak: the Mishnah is a legal code, not a chronicle, and its
silence is not surprising given its genre
(Section~\ref{sec:silence:mishnah}).

\item \textbf{The Dead Sea Scrolls community at Qumran}: deeply
engaged with messianic expectation and Jewish sectarianism.  No
mention.  The force is weak: the community's literary activity largely
predates the period of Jesus's supposed activity.
\end{enumerate}

The cumulative argument is strongest when limited to the first three
channels---sources that would plausibly mention a messianic figure
executed by the Roman prefect in their region and period---and weakest
for genre-inappropriate sources (the Mishnah) or chronologically
marginal ones (the Dead Sea Scrolls).  The argument does not depend on
any single weak channel.

The qualitative conclusion is defensible: the pattern of silence across
multiple traditions is \textit{more consistent with} the mythic-origin
model than with the historicist model, because the mythic-origin model
predicts silence while the historicist model must explain it through
auxiliary hypotheses about obscurity.  The cumulative weight of the
silence is real but difficult to quantify precisely.  It is one strand
of a cumulative case, not a standalone proof.

The argument from silence, properly formulated, is not a gap in the
evidence.  It is evidence.  Its cumulative weight favors the
mythic-origin hypothesis---but with less certainty than a maximalist
mythicist reading would claim.


%% ============================================================
\section{Prior Sources: The Jewish-Apocalyptic Matrix}
\label{sec:myth}
%% ============================================================

A central claim of this paper is that every narrative element of the
Jesus tradition can be traced to identifiable prior sources.  An older
popularization of this claim (associated with \citet{Frazer1890}) located
those sources primarily in pagan Mediterranean mystery religions---the
cults of Osiris, Attis, Adonis, Dionysus---and is what most readers
probably think of when they encounter the phrase ``mythicist hypothesis.''
That framing is the weakest version of the argument and is not the one
defended here.  As \citet{Fredriksen2017,Fredriksen2018} rightly
insists, Paul is a Second Temple Jewish apocalyptic thinker operating
within Jewish categories, not a syncretist blending in pagan myth.  The
Gospel authors, however they stood with respect to Jewish practice, are
demonstrably citing Jewish scripture as their primary source of
narrative material.  The honest argument is therefore not that
Christianity copied dying-and-rising mystery cults but that the Jesus
tradition is constructed from the scriptural, apocalyptic, and
mystical resources of Second Temple Judaism itself, with Hellenistic
literary conventions providing ambient narrative shaping secondary to
that primary Jewish matrix.

This section develops that argument.
Section~\ref{sec:myth:jewish} establishes the Jewish-apocalyptic matrix:
pre-existent heavenly mediators, angelomorphic Christology,
Philo's~\textit{Logos}, the Danielic Son of Man, and the Suffering
Servant.  Section~\ref{sec:myth:hb} catalogs the specific Hebrew Bible
templates from which Gospel narrative episodes are constructed.
Section~\ref{sec:myth:davidic} addresses the strongest specific
challenge to the mythic-origin reading of Paul---the Davidic-descent
claim in Romans~1:3---and engages \citet{Novenson2017} on Second Temple
messianic grammar.  Section~\ref{sec:myth:hellenistic} locates the
secondary Hellenistic resonances in their proper, supporting role.
Section~\ref{sec:myth:crucified} addresses the ``nobody would invent a
crucified Messiah'' argument.

%% ------------------------------------------------------------
\subsection{The Jewish-Apocalyptic Matrix as Primary Source}
\label{sec:myth:jewish}
%% ------------------------------------------------------------

The decisive observation of recent scholarship, from
\citet{Fredriksen2017} through \citet{Boyarin2012} and in the explicit
argumentative architecture of \citet{Carrier2014}, is that the
theological building-blocks of earliest Christian Christology are
already present in pre-Christian Second Temple Judaism.  Paul does not
reach outside his Jewish tradition for the raw materials of his
soteriology; he reaches deeper into it.

\paragraph{Pre-existent heavenly mediators.}
The idea of a divine or quasi-divine figure who exists alongside God,
functions as God's agent in creation and providence, and mediates
between the transcendent deity and human beings is extensively attested
in pre-Christian Jewish sources.  \citet{Boyarin2012} traces this
\enquote{binitarian} strand through the Danielic Son of Man
(Daniel~7:13--14), the \textit{Similitudes of Enoch} (1~Enoch~37--71),
the Melchizedek traditions at Qumran (11QMelch), the angel Metatron in
the Hekhalot tradition, and the Logos of Philo of Alexandria.  Paul's
Christ is not an exotic Hellenistic import into Judaism; he is a
recognizable member of the Jewish pre-existent-mediator family.  The
Philippians~2:5--11 hymn---\enquote{in the form of God\ldots emptied
himself\ldots taking the form of a servant\ldots therefore God highly
exalted him}---describes a descent-and-exaltation arc that maps onto
the Second Temple pattern of angelomorphic humiliation and vindication
(compare the Self-Glorification Hymn at Qumran, 4Q491c), not onto the
death-and-rebirth cycle of Attis.

\paragraph{Philo's Logos.}
\citet{Sterling2001} and the broader Philo commentary tradition
document Philo's Logos theology in detail.  The Logos is pre-existent,
is God's agent in creation, is called \enquote{son of God} and
\enquote{firstborn}, is the mediator through whom humans approach God,
and is sometimes described as \enquote{a second god}
(\textit{deuteros theos}, \textit{Questions on Genesis} 2.62).  Philo
wrote in Alexandria in the same decades as Paul's formation.  The
conceptual inventory from which early Christological language was
assembled was already Jewish-Hellenistic and available; no pagan
borrowing is required.

\paragraph{The Danielic Son of Man and the apocalyptic messiah.}
Daniel~7:13--14---\enquote{one like a son of man came with the clouds
of heaven\ldots and to him was given dominion and glory and kingdom}---
became a central messianic text in Second Temple apocalypticism.
The \textit{Similitudes of Enoch} (1~Enoch~37--71), dateable on
current best estimates to the late first century~BCE or early first
century~CE, identifies this Son of Man as a pre-existent, named, and
hidden heavenly figure who will appear at the end of days to execute
judgment.  This is not Christian material; it is the Jewish matrix
from which Christian material was drawn.

\paragraph{The Suffering Servant and vicarious death.}
Isaiah 52:13--53:12 describes a righteous figure who suffers vicariously
(\enquote{he was pierced for our transgressions, crushed for our
iniquities}), is \enquote{cut off from the land of the living,} and is
afterward vindicated.  Whether this passage was read messianically in
pre-Christian Judaism is contested, but the \textit{raw material} for a
theology of a righteous sufferer whose death atones for others is
already present in the Jewish canon.  \citet{Nickelsburg1972} showed
that Wisdom of Solomon~2--5 had independently developed a similar
pattern of persecuted-righteous-one vindication.

\paragraph{The significance of the matrix for mythic-origin
scenarios.}
The mythic-origin scenario does not require recourse to dying-and-
rising mystery cults.  Every element Paul claims about Christ---pre-
existence, descent, vicarious suffering, resurrection, exaltation,
cosmic lordship---is derivable from Jewish scriptural and apocalyptic
sources available in the first century.  This is also \textit{a
fortiori} the safer argument, since the Jewish sources are textually
secure (the Qumran material was effectively sealed by 70~CE and
shows no Christian influence) while the pagan parallels face genuine
scholarly contestation as to dating, comparability, and direction of
influence.

%% ------------------------------------------------------------
\subsection{The Revelatory Mechanism: Jewish Mystical Ascent and Pauline \textit{Apokalypsis}}
\label{sec:myth:mystical}
%% ------------------------------------------------------------

The mythic-origin scenario claims that Paul's Christ was known to the
earliest community through scriptural exegesis and revelation rather
than through historical memory.  This claim is commonly dismissed as
methodologically convenient---``revelation can explain anything''---but
the dismissal is not warranted.  Second Temple Judaism possessed a
specific, well-attested set of \textit{procedures} for generating
religious knowledge about heavenly realities.  These procedures have
been documented in detail by specialists on Jewish mysticism and
apocalypticism, though the literature has been surprisingly little
engaged by historical-Jesus scholars on either side.

\paragraph{Paul's own testimony on method.}
Paul explicitly and repeatedly tells us that his Christ-knowledge was
generated by two channels, revelation and scripture:
\enquote{I did not receive [the gospel] from any human source nor was I
taught it, but I received it through a revelation of Jesus Christ}
(\grk{di' apokalypse\=os I\=esou Christou}, Galatians~1:12).
\enquote{When God\ldots was pleased to reveal his Son in me
(\grk{apokalypsai ton huion autou en emoi}), so that I might proclaim
him among the Gentiles} (Galatians~1:15--16).
\enquote{I went up in response to a revelation} (\grk{kata
apokalypsin}, Galatians~2:2).
Paul places himself among \enquote{those who were apostles before me}
(Galatians~1:17), and the context throughout Galatians 1--2 frames
his knowledge as received directly rather than mediated through a
recently-living figure.

Paul also cites scripture as the source of his most compressed
Christological statements.  \enquote{Christ died for our sins
\textit{in accordance with the scriptures}\ldots and was raised on the
third day \textit{in accordance with the scriptures}}
(1~Corinthians~15:3--4).  The formula is apparently pre-Pauline; its
grammar is exegetical; scripture is explicitly the warranting source
of the content.

The combination of \textit{apokalypsis} with scriptural exegesis is
not an ad hoc pairing.  It is the attested epistemology of Second
Temple Jewish apocalyptic-mystical writing generally.

\paragraph{The Jewish mystical-ascent tradition.}
Pre-Christian and early-Christian-era Jewish literature documents a
specific set of practices and texts in which visionary ascent through
the heavens and exegetical decoding of scripture together generate
knowledge of heavenly realities that the community then articulates
and transmits.  The foundational modern study is \citet{Scholem1941},
supplemented by \citet{Gruenwald1980}, \citet{Rowland1982},
\citet{SchaferP2009}, and \citet{Elior2004}.

Representative texts and traditions:

\begin{itemize}[leftmargin=2em]
\item \textbf{1~Enoch} (portions pre-second-century BCE through
first century~CE), especially the \textit{Similitudes}
(1~Enoch~37--71): Enoch is taken up to heaven, shown cosmic secrets,
and identifies a named, hidden, pre-existent Son of Man who will
execute eschatological judgment.  This is not Christian material;
it is the Jewish matrix.
\item \textbf{The Apocalypse of Abraham} (first or second century
CE): Abraham undergoes angelic ascent, receives visions of the
cosmic order, and learns the identity of heavenly mediators.
\item \textbf{The Ascension of Isaiah} (first or early second
century; discussed in Section~\ref{sec:noncanon:ascension}): Isaiah
ascends through the seven heavens, sees the pre-existent Beloved,
and witnesses the Beloved's descent, crucifixion (in the sublunar
realm, in some readings), and return.  The text is formally a
mystical-ascent apocalypse, and some readers \citep{Carrier2014}
argue it preserves an earlier form of the Christian revelation
prior to its historicization.
\item \textbf{The Dead Sea Scrolls Self-Glorification Hymn}
(4Q491c, 4Q427): an otherwise-unnamed speaker claims to have
ascended to heaven, to sit enthroned among the angelic council,
and to have received divine knowledge.  The speaker's identity is
contested, but the \textit{procedure}---first-person mystical
ascent to the heavenly court, receiving knowledge and authority---
is exactly the procedure Paul describes for himself
(2~Corinthians~12:2--4: \enquote{I know a person in Christ who
fourteen years ago was caught up to the third heaven\ldots caught
up into Paradise and heard things that are not to be told}).
\item \textbf{The Merkabah and Hekhalot literature} (traditionally
dated late but with first-century-CE roots according to
\citealt{SchaferP2009} and \citealt{Elior2004}): detailed
instructions and narratives for mystical ascent to the
\textit{merkabah} (divine chariot-throne of Ezekiel~1), culminating
in vision of the seated figure on the throne and reception of
heavenly knowledge.
\item \textbf{Philo's own reports}: Philo describes ecstatic
experiences and visionary reception of philosophical-theological
truth (\textit{On the Creation}~71; \textit{Special Laws}~3.1--6;
\textit{On the Cherubim}~27), and his interpretive method for
scripture is allegorical decoding of a text understood as
encoding heavenly mysteries.
\end{itemize}

The common structure across these texts is: a visionary ascends (or
is taken up by an angelic guide) through celestial realms, receives
visions of heavenly figures and cosmic secrets, and is commissioned
to communicate those realities.  The resulting literature is
explicitly generated by these means.  Nothing in the procedure
requires or implies recent historical events; the content of the
vision is new knowledge, not recalled memory.

\paragraph{Paul's profile fits this matrix exactly.}
\citet{Segal1990} argued in detail that Paul's self-description
matches the profile of a Jewish mystical visionary: he reports
direct \textit{apokalypsis} of a heavenly Christ-figure
(Galatians~1:12--16), describes himself being caught up to the third
heaven where he hears unutterable things (2~Corinthians~12:2--4),
claims his knowledge was not mediated by human testimony
(Galatians~1:11--17), and frames his apostolic authority as derivative
from direct heavenly commissioning rather than from discipleship to
a teacher.  \citet{MortonSmith1978} argued further that the earliest
Jesus tradition itself bears the marks of a mystical-ecstatic
operational mode---an argument that has been controversial but that
remains serious in the scholarly literature.  \citet{Rowland1982}
situates Paul's self-presentation within the mainstream Jewish
apocalyptic-visionary tradition rather than as an isolated or exotic
phenomenon.

\paragraph{What this implies for historicity.}
The implication is not that Paul was lying or that his community
fabricated its claims.  It is that the \textit{epistemic mechanism}
by which Paul and his community came to know about their Christ is
the same epistemic mechanism by which the authors of 1~Enoch's
\textit{Similitudes}, the Self-Glorification Hymn, the Apocalypse of
Abraham, the Ascension of Isaiah, and the Merkabah literature came
to know about their heavenly figures.  None of those figures
required a recent historical instantiation for the revelatory
knowledge to be generated.  The earliest Christian community could
have come to know its Christ by the same route, and Paul
\textit{tells us} that that is how he came to know him.

The historicist response to this observation must be that Paul's
language about revelation is a theological gloss on what was in fact
community memory of a historical figure.  This is possible but
requires that Paul, who emphatically insists his knowledge came
\textit{not from human sources} (Galatians~1:11--12, 1:16--17), was
either lying, mistaken, or using technical theological language in a
sense he did not mean.  None of these is the most parsimonious
reading of the text.  The parsimonious reading is that Paul means
what he says about his epistemic source, that his procedure is the
well-attested Jewish mystical-apocalyptic procedure, and that the
content of his revelation is the same kind of content (a heavenly
mediator figure with cosmic soteriological functions) that the
comparable texts in his tradition also produce.

This does not prove that the Christ Paul revealed had no historical
instantiation.  It does establish that \textit{no historical
instantiation is required} for the evidentiary trail Paul leaves to
exist.  The historicity of Jesus is, in this framing, not a
necessary postulate to explain Paul's letters; it is a separate
question requiring independent evidence.

%% ------------------------------------------------------------
\subsection{Hebrew Bible Typology: The Specific Templates}
\label{sec:myth:hb}
%% ------------------------------------------------------------

The Jewish scriptural sources for the Jesus narrative have been discussed
in Section~\ref{sec:gospels:passion} in connection with the passion
narrative.  Here the point is generalized: the Jesus tradition is
constructed from Hebrew Bible typology not only in its death narrative
but across its entire arc.

The birth narratives draw on the births of Moses (Pharaoh's massacre /
Herod's massacre), Samuel (miraculous conception, dedication to God's
service), and Isaac (divine promise, elderly or virginal conception).
The ministry recapitulates Elijah-Elisha (Section~\ref{sec:gospels:mark}).
The entry into Jerusalem fulfills Zechariah~9:9.  The Last Supper draws
on Passover typology.  The crucifixion is built from Psalm~22 and
Isaiah~53.  The resurrection draws on Hosea~6:2 (``after two days he
will revive us; on the third day he will raise us up''), Jonah~1:17
(three days in the belly of the fish, explicitly cited in Matthew~12:40),
and the general pattern of divine vindication found throughout the
Psalms and prophetic literature.

The ascension draws on Elijah's ascent to heaven in a chariot of fire
(2~Kings~2:11) and Daniel's Son of Man ascending to the Ancient of Days
(Daniel~7:13--14).  The promise of return draws on Daniel~7, Zechariah~14,
and the general eschatological framework of late Second Temple Judaism.

When every major episode in the Jesus narrative can be derived from a
specific scriptural source, the hypothesis that historical events
\textit{preceded} the scriptural interpretation becomes less
parsimonious than the hypothesis that scriptural sources \textit{generated}
the narrative.  The tradition reads as a systematic program of
scriptural fulfillment---which is exactly how the Gospel authors
present it (Matthew's formula citations: ``this took place to fulfill
what had been spoken through the prophet\ldots'').  The question is
whether the authors were recognizing fulfillment in historical events or
constructing narrative to demonstrate fulfillment.  The evidence, as
analyzed throughout this paper, favors the latter.

%% ------------------------------------------------------------
\subsection{Davidic Messianism and Romans~1:3}
\label{sec:myth:davidic}
%% ------------------------------------------------------------

The single most difficult Pauline passage for the mythic-origin
scenario is Romans~1:3: Christ \enquote{came into being from the seed
of David according to the flesh} (\grk{tou genomenou ek spermatos
Dauid kata sarka}).  Davidic descent is a genealogical claim.
\citet{Novenson2012,Novenson2017} has documented systematically that
\enquote{messiah} language in Second Temple Judaism operates within
specific grammatical constraints: a messiah is a human figure, a
descendant of a specific royal line, anointed for specific offices.
Messianic language was not available as a detachable cosmic abstraction
in Paul's environment.  If Rom~1:3 is an authentic Pauline statement
about Christ, it anchors the Christ-figure in a specific human
genealogical line.

The mythicist response must be honest about the difficulty.  The
response runs along three lines.

\paragraph{Interpolation.}
Romans~1:3--4 is widely recognized as a pre-Pauline creedal fragment
quoted by Paul rather than composed by him (Section~\ref{sec:paul:dating};
\citealt{Jewett2007}).  Paul inherits this formula; he does not
originate it.  The formula belongs to the earliest Christian community,
but its phrasing may reflect that community's interpretive
commitments rather than Paul's own primary Christological categories.
Elsewhere in Paul, Davidic descent never functions argumentatively;
it is asserted once, in a formula, and not developed.

\paragraph{Celestial incarnation into Davidic flesh.}
\citet{Carrier2014} argues that \grk{kata sarka} (\enquote{according to
the flesh}) is a technical Pauline phrase that can locate an event in
the sublunar realm of flesh without committing to a specific earthly
biography.  A pre-existent celestial figure who takes on Davidic
lineage in the moment of descent into flesh is a coherent (if
counterintuitive) reading of the Jewish angelomorphic material catalogued
above.  This reading is not the most natural reading of Rom~1:3 in
isolation; it is the reading that harmonizes Rom~1:3 with the broader
pattern of Paul's discourse.

\paragraph{The remaining honest tension.}
Neither of these responses fully dissolves the difficulty.  Novenson's
point is not blunted by the creedal-fragment move: even if the formula
is pre-Pauline, Paul has endorsed it by including it in his prefatory
greeting to the Romans, and pre-Pauline creeds are evidence of the
\textit{earliest} stratum of Christian claim-making.  And the
celestial-incarnation reading of \grk{kata sarka} in Rom~1:3 is
defensible but strained.  This paper's conclusion acknowledges the
passage as one of the genuine challenges for mythicism.  A rigorous
mythic-origin scenario must either accept a weakened prior on
historicity grounds from Rom~1:3 or propose a defensible interpolation
hypothesis for the specific clause \grk{ek spermatos Dauid}.  We do
not pretend that this has been resolved.

A final observation: the difficulty of Rom~1:3 is shared with the
remaining Pauline tensions noted throughout this paper (Gal~4:4,
1~Cor~15:3--8, and some passages in 1~Corinthians).  These tensions
do not vanish under the mythicist scenario, but they are also not
neutral evidence for historicity.  They are specifically evidence that
Paul's earliest community had already moved---by the time Paul wrote in
the 50s---from a purely revelatory celestial Christ to one who had at
minimum been woven into a Jewish messianic-genealogical frame.  What
the evidence cannot distinguish is whether that move reflects
memory of a specific historical individual or reflects scriptural
exegesis identifying the heavenly mediator as, necessarily, of David's
line (Jer~23:5, 33:15; Isa~11:1--10; Ps~89; 2~Sam~7).

%% ------------------------------------------------------------
\subsection{Hellenistic Resonances as Secondary Context}
\label{sec:myth:hellenistic}
%% ------------------------------------------------------------

A complete account of the ambient literary and religious environment in
which the Jesus tradition took shape includes, secondarily, a set of
Hellenistic conventions and religious patterns.  These are not the
primary source of the tradition---the Jewish matrix of
Section~\ref{sec:myth:jewish} is---but they represent the shared
narrative grammar of the Mediterranean world within which the Jesus
tradition was articulated, circulated, and received.

\paragraph{The \textit{theios an\=er} biographical template.}
Greco-Roman literary culture had a well-established pattern for
biographies of divinely-endowed figures (\textit{theios an\=er},
\enquote{divine man})---miraculous birth, wisdom teaching,
miracle-working, conflict with authorities, noble death, and
post-mortem appearances or ascent \citep{Bieler1935,Tiede1972}.
Apollonius of Tyana, Pythagoras, and Alexander the Great attracted
biographical traditions conforming to this pattern.  Philostratus's
\textit{Life of Apollonius} (\circa 220~CE, drawing on earlier
material) is the closest extant literary analogue to the Gospels, and
the structural parallels are extensive.  The Gospel authors were not
writing in a cultural vacuum; they were writing in a literary world
where biographical narratives of divine-human figures had conventional
shapes.  The Jesus story conforms to those shapes in every major
element.  This does not prove the subject was fictional---historical
subjects can be described in conventional terms---but it removes any
presumption that the narrative's biographical specificity reflects
specifically historical (as opposed to conventionally literary)
content.

\paragraph{The dying-and-rising-god pattern (carefully qualified).}
A generation of mythicist writing prominently argued, following
\citet{Frazer1890}, that Christianity drew centrally on a Mediterranean
pattern of deities who die and return to life: Osiris, Attis, Adonis,
Dionysus, Tammuz/Dumuzi.  This argument has been heavily criticized.
\citet{Smith2005}, in the most widely cited recent critique, argued
that the \enquote{dying-and-rising god} category was substantially a
construction of nineteenth- and twentieth-century scholarship that
over-read the sources.  \citet{Mettinger2001} subsequently defended a
more limited version of the category---genuine pre-Christian
dying-and-rising deities in Mesopotamian (Dumuzi), Ugaritic (Baal),
and Phoenician (Melqart) sources, with the Adonis case probably
genuine---while agreeing that the older, broader Frazerian
generalization was indefensible.

The present paper takes the cautious position that some limited
Mediterranean dying-and-rising tradition is real but is \textit{not}
the primary source of Pauline or Gospel material.  Paul's actual
Christology, as developed in Section~\ref{sec:myth:jewish}, is
assembled from Jewish apocalyptic-mystical resources.  The Hellenistic
dying-and-rising pattern may have made Paul's gospel culturally
intelligible to gentile audiences (this is a plausible reading of
1~Cor~1:23, \enquote{Christ crucified\ldots foolishness to the
Gentiles}), but it does not supply the generative theological machinery.
Historicist authors who treat \enquote{parallels with Attis and Osiris}
as the central mythicist claim are attacking a weaker version of the
argument than the one this paper defends.  \citet{Fredriksen2017} is
correct that Paul is not a syncretist drawing on pagan cult; the
mythicist case does not require that he was.

%% ------------------------------------------------------------
\subsection{The ``Nobody Would Invent a Crucified Messiah'' Argument}
\label{sec:myth:crucified}
%% ------------------------------------------------------------

The argument from the ``scandal'' of the cross is perhaps the most
influential single argument for historicity.  It runs as follows: Jewish
messianic expectation anticipated a triumphant royal or priestly figure,
not a crucified criminal.  Crucifixion was the most shameful form of
Roman execution.  No one within a Jewish context would invent a Messiah
who was crucified; therefore, the crucifixion must be historical, and
the early Christian claim that the crucified Jesus was the Messiah must
have arisen from the need to make sense of an actual event
\citep{Hengel1977,Wright1996}.

This argument has been treated as nearly decisive for decades.  It has
serious problems.

First, it assumes complete knowledge of Second Temple Jewish sectarian
thought.  The Dead Sea Scrolls, discovered beginning in 1947, revealed
entire theological traditions previously unknown to scholarship.  The
claim that ``no Jew would expect a dying Messiah'' is an argument from
ignorance of a tradition that was far more diverse than our evidence
allows us to map exhaustively.

Second, the argument ignores the Suffering Servant tradition of
Isaiah~52:13--53:12, which was available in the Jewish canon and
describes a righteous figure who suffers vicariously, is despised and
rejected, is ``cut off from the land of the living,'' and is afterward
vindicated by God.  Whether or not this passage was interpreted
messianically in pre-Christian Judaism (a debated question), it
provided the raw material for a theology of vicarious suffering and
vindication that could be applied to a messianic figure by anyone doing
creative scriptural exegesis---which is precisely what the earliest
Christians demonstrably did.

Third, the argument assumes an exclusively Jewish audience and an
exclusively Jewish frame of reference.  But Paul---our earliest
source---is writing to mixed communities across the Mediterranean world,
communities in which dying-and-rising gods were familiar and religiously
compelling.  In this broader context, a dying and rising savior was
recognizable.  However, the historicist has a genuine counter-point
here: crucifixion \textit{specifically}---as distinct from divine death
generally---carried a particular Roman stigma of degradation and
criminality that voluntary divine death (Osiris, Attis) or heroic
dismemberment (Dionysus) did not.  Paul's own statement that the
crucified Christ is ``a stumbling block to Jews and foolishness to
Gentiles'' (1~Corinthians~1:23) plausibly reflects real audience
resistance, not merely rhetorical preemption.  The mythicist must
explain why the tradition chose crucifixion specifically rather than
some other form of death.

The strongest mythicist answer centers on Deuteronomy~21:23: ``anyone
hung on a tree is under God's curse.''  Paul explicitly cites this
passage in Galatians~3:13: ``Christ redeemed us from the curse of the
law by becoming a curse for us---for it is written, `Cursed is everyone
who hangs on a tree.'\,''  This is not an afterthought; it is the
\textit{engine} of Paul's soteriology.  The entire logic of vicarious
atonement depends on the savior bearing a curse on behalf of others.
For this theology to work, the mode of death must be one that is
\textit{scripturally cursed}.  Stoning, burning, beheading---none of
these carry the specific scriptural curse that hanging on a tree does.
Crucifixion was not chosen despite being embarrassing; it was chosen
\textit{because} it was the one mode of death that could be scripturally
identified as ``cursed by God,'' making it the necessary vehicle for a
theology of substitutionary curse-bearing.  On this reading, the
``scandal'' of the cross is not evidence of historicity but evidence of
theological engineering: the tradition selected the mode of death that
generated the maximum soteriological power from available scriptural
resources.

Additionally, crucifixion was the most visible, public, and humiliating
form of Roman execution---designed not merely to kill but to display
power and inflict shame.  A savior who died in obscurity would lack
the dramatic force that the tradition required.  The extremity of the
degradation is theologically generative: it maximizes the paradox of
divine humiliation and exaltation that structures the Christ hymn in
Philippians~2:5--11 (``he humbled himself and became obedient to the
point of death---even death on a cross''), where ``even death on a
cross'' functions as a climactic intensifier, not as a reluctant
historical admission.

Fourth, the argument has a hidden premise: that embarrassment implies
historicity.  This is not generally true.  Mythological traditions
routinely include elements that are ``embarrassing'' by the standards of
some audiences while being powerful and attractive to others.  The death
of Osiris is violent and degrading; the self-castration of Attis is
shocking; the dismemberment of Dionysus is grotesque.  These elements
were not historical concessions; they were mythologically generative.
The ``scandal'' of the cross may function identically: a narratively
powerful element that creates theological meaning precisely through its
extremity.  The parallel is imperfect---crucifixion was a specific,
contemporary, politically charged mode of execution---but the principle
holds: transgressive and shocking elements in mythological traditions
do not require historical events to explain their presence.

The ``nobody would invent a crucified Messiah'' argument is weaker than
the consensus maintains.  The Deuteronomy~21:23 argument provides a
specific, text-grounded reason why crucifixion was selected, and it is
Paul himself who makes the connection explicit.  However, the argument
retains some force: the specificity of crucifixion as a Roman punishment,
with its particular political associations, is not fully explained by
the scriptural-curse theory alone.  This is a point of genuine tension
rather than a decisive victory for either side.


%% ============================================================
\section{Source Trajectory: From Mythic to Historical}
\label{sec:trajectory}
%% ============================================================

%% ------------------------------------------------------------
\subsection{The Direction of Development}
\label{sec:trajectory:direction}
%% ------------------------------------------------------------

A fundamental prediction distinguishes the two hypotheses.  If a
historical Jesus existed and was subsequently mythologized, we would
expect the earliest sources to be the most historically grounded and
the latest sources to show increasing mythological embellishment.  If
a mythic figure was subsequently historicized, we would expect the
earliest sources to be the most mythic and the latest sources to show
increasing historical specificity.

The evidence matches the second prediction, not the first.

\begin{itemize}[leftmargin=2em]
\item \textbf{Paul} (50s~CE): the earliest source.  Jesus is a cosmic
figure known through revelation and scripture.  No birthplace, no
ministry, no named contemporaries, no narrative of earthly life.  The
most mythic, least historical layer.

\item \textbf{Mark} (\circa 70~CE): the first narrative.  A minimal
biographical framework is introduced---Galilee, Jerusalem, crucifixion
under Pilate---but much of the content is constructed from scriptural
templates.  No birth narrative, no resurrection appearances (the
original ending at Mark~16:8), no genealogy.

\item \textbf{Matthew and Luke} (\circa 80--90~CE): birth narratives
are added (mutually contradictory, constructed from proof-texts).
Genealogies are supplied.  Post-resurrection appearances are narrated.
The historical framework expands significantly.

\item \textbf{John} (\circa 90--100~CE): a different chronology is
introduced (three-year ministry vs.\ one year), new episodes are added
(Lazarus, the Samaritan woman, the wedding at Cana), and extended
theological discourses are placed in Jesus's mouth.

\item \textbf{Later apocrypha} (second century and beyond): childhood
miracles (Infancy Gospel of Thomas), extended dialogues with Pilate
(Acts of Pilate), detailed accounts of the descent into hell (Gospel
of Nicodemus), family details, and biographical elaboration continue
to accumulate.
\end{itemize}

The trajectory is unmistakable: from cosmic and mythic to biographical
and historically situated, from abstract to concrete, from silence about
earthly details to proliferating specification.  This is the signature
of euhemerization---the process by which a mythic figure is
progressively anchored in historical time and place---not of
mythologization, which would produce the reverse trajectory.

The historicist response is that the later elaboration represents
legendary accretion around a historical core, not euhemerization.  But
this response assumes the conclusion: it presupposes that the earliest
layer (Paul) contains a historical figure who is simply not discussed
in biographical terms, rather than a mythic figure who has not yet
been historicized.  The direction of the trajectory cannot distinguish
between these interpretations \textit{by itself}, but in combination
with the evidence analyzed in previous sections---the celestial
character of Paul's Christ, the scriptural construction of the Gospels,
the silence of non-Christian sources---the euhemerization model provides
a more coherent and parsimonious account.

%% ------------------------------------------------------------
\subsection{Euhemerization as a Documented Process}
\label{sec:trajectory:euhemerization}
%% ------------------------------------------------------------

Euhemerization---named after the Greek mythographer Euhemerus
(\circa 300~BCE), who argued that the gods were originally human kings
deified after death---is a well-attested process in the ancient world.
The reverse process---the historicization of mythic figures---is also
documented, though less commonly discussed.

Romulus, the legendary founder of Rome, was originally a mythic figure
(son of Mars, suckled by a she-wolf, ascended to heaven as the god
Quirinus) who was progressively historicized into a ``real'' king with
specific dates, political actions, and a detailed biography.  No
serious historian treats Romulus as historical, despite the extensive
biographical tradition.  The parallel with Jesus illustrates the
principle: elaborate biographical detail does not require a historical
person; it requires a mythic tradition with cultural motivation to
anchor itself in history.

However, the Romulus parallel has a significant disanalogy that must be
acknowledged: the historicization of Romulus took centuries, while the
Jesus tradition produced a biographical narrative (Mark) within
approximately forty years of the supposed events, and Paul's letters---
which already presuppose established communities worshipping
Christ---appear within twenty to twenty-five years.  This is fast for
euhemerization, and the speed of development is a genuine challenge for
the mythicist model.  The historicist argues that the rapid emergence of
a biographical tradition is more easily explained by the existence of a
historical person around whom the tradition formed.

Several considerations mitigate, though they do not eliminate, this
challenge.

First, the relevant time-compression is narrower than it appears.
Paul's letters (50s~CE) do \textit{not} show historicization: Paul's
Jesus is cosmic, not biographical.  The historicization occurs between
Paul and Mark---a window of perhaps fifteen to twenty years.  The
question is whether a mythic figure can be placed in a specific
historical setting within one generation.  This is fast, but not
unprecedented.

Second, the Cargo Cults of Melanesia, documented by anthropologists in
the twentieth century, provide direct evidence that rapid
historicization is possible.  In several cases, mythic or divine figures
were assigned specific names, biographical details, and historical
contexts within a single generation, producing traditions that look,
from the inside, like historical memory but are demonstrably mythic in
origin \citep{Worsley1957,Lindstrom1993}.  The conditions that
facilitated this---cultural upheaval caused by contact with a dominant
foreign power, intense eschatological expectation, communities under
existential stress---arguably obtained in first-century Roman Palestine.

Third, the specific conditions of the early Christian movement were
unusually conducive to rapid historicization.  The movement was
geographically dispersed across the Mediterranean within two decades
(Paul's letters attest this).  It encountered diverse audiences who
valued historical narrative.  It operated within a scriptural
tradition that provided ready-made biographical templates (the
Elijah-Elisha cycle, the Suffering Servant, the Psalmic righteous
sufferer).  And it faced apologetic pressure: a movement claiming a
cosmic Lord who had defeated death would face the question ``when and
where did this happen?'' from every audience it addressed.  The Gospel
of Mark can be understood as the first systematic answer to that
question---a biographical narrative assembled from scriptural templates
to give the cosmic Christ a historical address.

Fourth, the relevant comparison is not with Romulus (where
historicization occurred in a literate culture with stable institutions)
but with new religious movements in periods of rapid social change.
Joseph Smith produced a detailed historical narrative for the Book of
Mormon within a few years.  The Sabbatean movement reinterpreted
Sabbatai Zevi's apostasy into a mythological framework within months.
Religious movements under pressure can produce elaborate narrative
constructions very quickly.

These considerations do not fully resolve the time-compression problem.
The speed of the Paul-to-Mark transition remains a point where the
historicist model has a structural advantage: a historical person
provides an obvious catalyst for rapid narrative development that the
mythicist model must explain through other mechanisms.  But the
mechanisms are available, and the speed is not so extreme as to make
the mythicist model implausible---merely more demanding of explanation.

The mechanism by which historicization occurs is straightforward.  A
community that worships a mythic figure has practical reasons to anchor
that figure in historical time and place: it makes the figure more
relatable, more authoritative, and more resistant to charges of
fabrication.  As the movement grows and encounters audiences that value
historical narrative (as Roman culture did), the pressure to
historicize increases.  The Gospel of Mark can be understood as the
first major expression of this pressure within the Christian movement:
the creation of a biographical narrative for a figure who had previously
been known only through revelation and scripture.  But whether the
twenty-to-forty-year window is sufficient for this process---without a
historical catalyst---remains an open question that the mythicist model
has not fully resolved.


%% ============================================================
\section{The Sociology of the Consensus}
\label{sec:sociology}
%% ============================================================

If the evidence for the historical Jesus is as weak as this paper
contends, an explanation is required for the near-universal consensus
in its favor.  The explanation lies not in the evidence but in the
institutional and sociological structure of the field.

%% ------------------------------------------------------------
\subsection{Theological Origins of the Discipline}
\label{sec:sociology:origins}
%% ------------------------------------------------------------

The academic study of the ``historical Jesus'' originated in and
has been conducted primarily within institutions with theological
commitments.  The first and second ``quests'' for the historical Jesus
(\citealt{Schweitzer1906}; \citealt{Kasemann1954}) were conducted
almost entirely by Protestant theologians working within divinity
schools and seminary faculties.  The third quest, while broader in
institutional base, remains heavily centered in departments of religious
studies and theological faculties where Christianity is not merely an
object of study but the institutional \textit{raison d'\^etre}.

This matters because the historicity of Jesus is not merely an academic
question in these settings; it is a presupposition of the enterprise.
A department of historical Jesus studies that concluded there was no
historical Jesus would have undermined its own reason for existence.
This is not a conspiracy; it is an ordinary institutional dynamic in
which certain conclusions are structurally discouraged by the conditions
of employment, funding, and professional advancement.

\citet{Thompson2005}, a distinguished Old Testament scholar at the
University of Copenhagen, reported that his questioning of the
historicity of the patriarchs and of Jesus provoked professional
hostility disproportionate to the evidential merits of his arguments.
\citet{Brodie2012}, a Dominican priest and New Testament scholar,
described being effectively silenced within his order after publishing
his mythicist conclusions in \textit{Beyond the Quest for the Historical
Jesus}.  These are not isolated anecdotes; they illustrate the
professional costs of dissent on this particular question.

%% ------------------------------------------------------------
\subsection{Paradigm Maintenance and Institutional Pressure}
\label{sec:sociology:paradigm}
%% ------------------------------------------------------------

\citet{Kuhn1962} described the process by which scientific paradigms
are maintained: anomalies are accommodated through auxiliary hypotheses,
fundamental assumptions become invisible because they are absorbed
during training rather than independently evaluated, and challenges to
the paradigm are treated as challenges to the competence or good faith
of the challenger rather than as evidential contributions.

The historicity consensus exhibits all of these features.  Anomalies---the
silence of Paul on biographical details, the absence of non-Christian
attestation, the scriptural construction of the narrative---are
accommodated by auxiliary hypotheses (``Paul's letters are occasional
correspondence,'' ``Jesus was too obscure,'' ``the authors used
scripture to describe real events'').  Each auxiliary hypothesis is
individually possible, but their accumulation constitutes exactly the
pattern of paradigm maintenance that Kuhn described: the core
assumption is protected by an expanding shell of \textit{ad hoc}
defenses.

The treatment of mythicists exemplifies the sociological dynamics of
paradigm protection.  \citet{Ehrman2012} dismisses mythicism as a
position held by no ``competent scholar,'' defining competence in terms
so narrow as to exclude anyone who disagrees.  He repeatedly invokes the
consensus as evidence, a maneuver that would be recognized as circular
in any other field.  The rhetoric is not that of a scholar engaging with
evidence; it is that of a paradigm defender policing the boundaries of
acceptable opinion.

This is not unique to Ehrman.  The professional response to
\citet{Carrier2014}---a peer-reviewed monograph published by an academic
press, written by a scholar with a doctorate in ancient history from
Columbia University---has been largely dismissal rather than engagement.
The book has received few formal academic reviews relative to its
significance.  This is the sociological signature of a claim that the
field does not wish to address on its merits.

%% ------------------------------------------------------------
\subsection{Engagement with Anti-Mythicist Scholarship}
\label{sec:sociology:engagement}
%% ------------------------------------------------------------

Intellectual honesty requires acknowledging the most sophisticated
historicist responses to mythicism rather than caricaturing them.
Four bodies of work deserve specific notice.

\citet{Ehrman2012} remains the standard broadly-accessible academic
rebuttal.  Its strengths are its clarity and its marshaling of the
source-plurality and minimal-facts arguments.  This paper engages its
key moves in §\ref{sec:paul:brother} (the \grk{adelphos tou kyriou}
passage), §\ref{sec:gospels:q} (the independence of the supposed
sources), and §\ref{sec:myth:crucified} (the crucifixion-as-scandal
argument).  Its reliance on consensus as evidence in its own right
has been discussed in §\ref{sec:method:authority}.  Where Ehrman's
specific arguments survive scrutiny---notably the cumulative structure
of Paul's references to ``the twelve'' and ``those who were apostles
before me'' (Gal~1:17, 1~Cor~15:5--7), and the evidentiary force of
Paul's acknowledged meeting with Peter and James in Jerusalem (Gal~1--2,
conducted twenty-odd years after any supposed crucifixion and
constituting a chain of acquaintance at minimum two persons removed
from Jesus)---we concede the force and engage the substance.  Where
the arguments reduce to consensus-claims, we note that reduction.

\citet{Casey2014} offers the most combative historicist critique of
mythicism.  Casey argues that Aramaic sources underlie the Greek
Gospels, that the tradition preserves genuine memories transmitted in
Aramaic-speaking communities, and that mythicists lack the linguistic
competence to evaluate the evidence.  The Aramaic-source argument has
real force: it suggests that at least some Gospel material passed
through an Aramaic-speaking milieu, which is a non-trivial constraint
on any mythic-origin scenario.  It does not, however, establish
historicity on its own: Aramaic transmission is consistent with an
Aramaic-speaking sectarian community developing its theology through
scriptural exegesis.  Casey's \textit{ad hominem} attacks on mythicist
credentials do not address the substance of the arguments.

\citet{Meier1991} (and its subsequent volumes) represents the most
comprehensive historicist reconstruction to date.  Meier's methodological
caution is admirable, and his willingness to declare large portions of
the tradition unrecoverable is more honest than many historicist
treatments.  Meier's project, however, presupposes historicity rather
than establishing it: he asks \enquote{what can we know about the
historical Jesus?} rather than \enquote{did a historical Jesus exist?}
That framing is defensible within a research programme that has
accepted historicity as a starting premise; it does not address the
question this paper raises about whether the premise is warranted.

\citet{Allison2010} and \citet{Allison2021} are, in the judgement of
this paper, the strongest sustained historicist case currently in
print.  Allison is as methodologically sceptical about the criteria
of authenticity as any mythicist (§\ref{sec:method:criteria}) and has
explicitly moved toward memory-theory as the post-criteria framework
within which historicist reconstruction should proceed.  Allison's
claim is that the \textit{gestalt} of the tradition---its recurring
motifs, its stable narrative shape across otherwise-independent
witnesses---is better explained by the memory of a historical figure
than by purely literary construction.  This paper grants that
gestalt-level arguments are the most sophisticated form of the
historicist case and that they cannot be dismissed by pointing to
individual episodes' literary derivation.  We contend, however, that
the same gestalt features can be explained by a dense scriptural-
and-revelatory generative process operating in a cohesive sectarian
community, and that Allison's argument requires the additional
assumption that memory rather than exegesis is the more parsimonious
source of the gestalt---an assumption that requires independent
defence, which is the argument of §\ref{sec:myth}.

Finally, it is worth naming the Kuhnian dynamic that sustains
consensus on unusually-embedded premises \citep{Kuhn1962}: the
premises least subject to critical re-examination are not those least
well supported but those most deeply absorbed during disciplinary
formation, because they are the premises on which both believers and
sceptics within a given field have historically agreed.  The
historicity of Jesus has that status: it is the shared starting
premise of confessional scholarship, secular scholarship, and even
much critical scholarship since Schweitzer.  Premises that hold that
position are unusually resistant to scrutiny.  This is not an
indictment of any individual scholar's competence or good faith; it
is an observation about how any discipline handles its most deeply
embedded assumptions.  It is not evidence against historicity.  It
is the reason this paper insists on re-examining the primary
evidence directly rather than relying on the weight of consensus to
settle the question.


%% ============================================================
\section{A Constructive Alternative: The Mythicist Origin Scenario}
\label{sec:origins}
%% ============================================================

A responsible challenge to the historicist consensus must do more than
identify weaknesses in the evidence; it must offer a constructive
alternative that explains the origin of Christianity without a
historical founder.  The following sketch draws on the work of
\citet{Doherty1999,Doherty2009}, \citet{Carrier2014}, and earlier
scholars in the tradition, while noting points of genuine uncertainty.

\paragraph{Stage 1: The celestial Christ (\circa 30s--50s~CE).}
A small Jewish sectarian community, steeped in Second Temple scriptural
exegesis, developed a soteriology centered on a celestial intermediary
figure.  This figure was assembled from existing Jewish theological raw
materials: the Suffering Servant of Isaiah~52--53, the vindicated
righteous one of Wisdom~2--5, the Danielic Son of Man (Daniel~7:13--14),
the Logos/Wisdom tradition developed by Philo and earlier Jewish
thinkers, and angelomorphic figures such as those in 1~Enoch and the
Dead Sea Scrolls.  The community understood this figure---named
\grk{I\=esous} (Joshua/Jesus, meaning ``Yahweh saves'')---as a
pre-existent divine being who descended into the domain of flesh, died
as a vicarious sacrifice ``in accordance with the scriptures,'' was
raised by God, and was exalted to cosmic lordship.

The community's knowledge of this figure came through two channels:
scriptural exegesis (reading the Hebrew Bible as encoding the Christ
mystery) and visionary revelation (ecstatic experiences interpreted as
encounters with the risen Lord).  Paul's letters are the primary
evidence for this stage: they describe exactly this kind of figure,
known through exactly these channels.

\paragraph{Stage 2: Community formation and diversification
(\circa 40s--60s~CE).}
The movement spread across the Mediterranean, establishing communities
in major cities (Corinth, Thessalonica, Philippi, Rome).  Different
communities developed different emphases---some more
apocalyptic, some more wisdom-oriented, some more Hellenized---but all
shared the core soteriology of the dying and rising Christ.  Leadership
disputes emerged (Paul vs.\ the Jerusalem ``pillars''), ritual practices
diversified (the Lord's Supper, baptism, prophecy), and the tradition
began generating ``sayings of the Lord'' through prophetic revelation
attributed to the risen Christ.  This accounts for the small number of
``commands of the Lord'' Paul cites (1~Corinthians~7:10, 9:14), as well
as the wider body of sayings that would eventually be compiled in Q or
its equivalent.

\paragraph{Stage 3: Historicization in Mark (\circa 70~CE).}
The destruction of the Jerusalem Temple in 70~CE was a theological
crisis for Judaism and for Jewish-Christian communities.  In this
context, the author of Mark composed the first narrative gospel: a
biographical account of the Christ figure, set in recent Palestinian
history, structured around a journey from Galilee to Jerusalem
culminating in crucifixion under Pontius Pilate.  Mark built his
narrative from the available materials: the Elijah-Elisha cycle for the
ministry, the Psalms and Isaiah for the passion, and the theological
traditions of the community for the overall framework.  Mark's Jesus
\textit{is} the celestial Christ of Paul, given a historical address.

The motivations for historicization were practical: a narrative set in
identifiable times and places was more compelling to audiences
accustomed to Greco-Roman biography; it answered the question ``when and
where did this happen?''; it provided a framework for organizing the
scattered sayings and teachings attributed to the Lord; and it made the
tradition more resistant to the charge that the Christ was merely a
philosophical abstraction.

\paragraph{Stage 4: Elaboration (\circa 80--100~CE and beyond).}
Matthew and Luke independently expanded Mark's narrative, adding birth
narratives (from different scriptural templates), genealogies,
resurrection appearances, and additional teaching material.  John
produced a theologically distinct narrative with a different chronology
and new episodes.  Subsequent centuries produced further elaboration:
childhood miracles, detailed Pilate dialogues, harrowing-of-hell
narratives, and the full apparatus of Christian hagiography.

\paragraph{Assessment.}
This scenario is \textit{constructive}: it provides a coherent account
of how Christianity could have originated and developed without a
historical founder.  It is \textit{consistent} with all the evidence
analyzed in this paper: Paul's celestial Christ, the scriptural
construction of the Gospels, the silence of non-Christian sources, the
mythological parallels, and the source trajectory from mythic to
historical.  It is \textit{not proven}: alternative scenarios (including
historicity) remain possible.  But it demonstrates that the mythicist
hypothesis is not merely a negation of historicism; it is a positive
historical model with explanatory power.

The scenario's weakest points---the speed of the Paul-to-Mark
transition (Section~\ref{sec:trajectory:euhemerization}), the
specificity of the 1~Corinthians~15 witness list
(Section~\ref{sec:paul:character}), and the two ``commands of the Lord''
(Section~\ref{sec:paul:silence})---have been openly acknowledged
throughout this paper.  These remain genuine challenges.  But they are
challenges of detail, not of structure: the overall architecture of the
mythicist scenario fits the evidence, and the individual difficulties
admit possible (if not yet fully satisfying) explanations.


%% ============================================================
\section{Conclusion}
\label{sec:conclusion}
%% ============================================================

The evidence for the historical existence of Jesus of Nazareth has been
examined across every available category: the Pauline epistles, the
canonical Gospels, non-canonical Christian sources, non-Christian
attestation, and the documentary traditions that should contain
references but do not.  The findings may be summarized as follows.

The Pauline corpus, the earliest stratum of Christian writing, primarily
describes a celestial, revelatory figure known through scripture and
visionary experience.  The single passage most frequently cited as
decisive for historicity---Galatians~1:19, ``brother of the Lord''---is
more ambiguous than the consensus acknowledges, particularly in light of
the plural ``brothers of the Lord'' in 1~Corinthians~9:5.  However,
several Pauline passages pose genuine challenges for the mythicist
model: Galatians~4:4 (``born of a woman, born under the law''),
Romans~1:3 (``descended from David according to the flesh''), and the
specificity of the witness list in 1~Corinthians~15:3--8.  Mythicist
readings of these passages are possible but require interpretive work
that the historicist readings do not, and intellectual honesty requires
acknowledging this asymmetry.

The Gospel narratives are demonstrably constructed from Hebrew Bible
templates.  The passion narrative maps onto Psalm~22, Isaiah~53,
Wisdom~2--5, and Zechariah point by point.  The birth narratives are
mutually contradictory and built from proof-texts.  The ministry
episodes recapitulate the Elijah-Elisha cycle.  This templating does
not by itself prove the non-historicity of a subject---historical
persons can be described through scriptural grammars---but it does
establish that the Gospels are not independent historical testimony.
They testify to the theological constructions of late-first-century
communities, not to witness memory of an event.  The historicist
evidential burden therefore cannot be carried by the Gospels; it must
be carried by Paul and the non-Christian attestation, each of which
is constrained and limited as analysed above.

The non-canonical Christian sources---Thomas, the Didache, 1~Clement,
Barnabas---reinforce this finding.  The earliest non-Pauline sources
construct their Jesus from scripture and creed, not from biographical
memory.

The non-Christian attestation is weaker than the consensus typically
claims but not empty.  The \textit{Testimonium Flavianum} is partially
or wholly interpolated.  The James passage in Josephus is accepted by
most textual scholars as substantially genuine and constitutes the
strongest non-Pauline evidence for historicity, though its evidential
value is limited.  Tacitus almost certainly derives his information from
Christian tradition.  Collectively, the non-Christian sources confirm
that by the late first and early second centuries the historicized
narrative was in wide circulation---a datum that constrains both
hypotheses.

The silence of contemporaneous sources---Philo, Justus of Tiberias, the
Mishnaic tradition, Roman administrative records---constitutes
disconfirming evidence whose qualitative weight favors the mythic-origin
hypothesis.  However, the individual silences vary in strength: some
(Justus of Tiberias) are suggestive, while others (the Mishnah, a legal
code rather than a chronicle) are weak.  The cumulative argument from
silence retains force but cannot be honestly quantified with the
precision that a naive Bayesian calculation might suggest.

Every element of Paul's Christology and of the Gospel narrative has
identifiable prior sources in pre-Christian Second Temple Jewish
literature: the Danielic Son of Man, the Suffering Servant of Isaiah,
the Wisdom-pattern of the persecuted-and-vindicated righteous one,
angelomorphic mediator traditions, Philo's Logos, Hebrew Bible
typology across the narrative arc, and Davidic messianic grammar.
The Hellenistic literary environment supplied the ambient biographical
conventions within which those Jewish materials were articulated, but
the generative theological machinery is Jewish, not pagan.  The
earlier mythicist framing that located the primary source in pagan
dying-and-rising mystery cults is the weaker form of the argument and
is not the one this paper defends.  The specific mode of death---Roman
crucifixion---poses a genuine question for the mythic-origin model, but
is plausibly explained by Deuteronomy~21:23 as exploited in
Galatians~3:13: a death specifically under scriptural curse is the
mode of death required for a substitutionary-curse theology.

The source trajectory runs from mythic to historical---from Paul's
cosmic Christ through Mark's minimal narrative to the elaborated
biographies of Matthew, Luke, and John---which is consistent with
euhemerization.  However, the speed of this development (a biographical
narrative within forty years) is fast for euhemerization and constitutes
a genuine challenge that the mythicist model must address.

The scholarly consensus in favor of historicity is sustained partly by
the institutional structure of the field---its theological origins, its
paradigmatic assumptions, and the professional costs of dissent.  This
sociological observation does not make the consensus false, but it means
the consensus cannot be invoked as independent evidence for historicity.

\vspace{1em}

\textbf{Overall assessment.}  This paper has argued that the evidence
for the historical Jesus is substantially weaker than the consensus
claims, and that the mythic-origin model provides a parsimonious account
of much of the evidence---particularly the Pauline material, the
scriptural construction of the Gospels, and the silence of
contemporaneous sources.

It has not argued, and does not claim, that mythicism has been
\textit{proven} more probable than historicity.  Several features of the
evidence---the Pauline passages that admit historical readings
(Galatians~4:4, Romans~1:3, 1~Corinthians~15:3--8), the specificity of
crucifixion, the speed of the tradition's development, and the Josephan
James passage---pose genuine difficulties for the mythicist model that
have not been fully resolved.

The responsible conclusion is that the evidence does not establish the
historicity of Jesus with the confidence the consensus claims.  The
mythicist hypothesis is a serious alternative that the evidence does not
rule out.  Responsible scholarship must acknowledge that the historicity
question is legitimately disputed and that the mythicist model---while
not proven---provides a coherent, evidence-based account of Christian
origins that merits engagement on its merits rather than dismissal by
appeal to authority.

To be explicit about what this paper has \textit{not} argued: it has
not argued that Jesus definitely did not exist.  It has not argued that
mythicism is proven.  It has not argued that all historicist scholars
are incompetent or acting in bad faith.  It has argued that the evidence
is insufficient to support the near-certainty the consensus projects,
that the mythicist hypothesis is a serious alternative that survives
scrutiny of its weaknesses, and that the scholarly consensus is
sustained partly by institutional and methodological inertia rather than
by evidential weight alone.

The historicity of Jesus is not established beyond reasonable historical
doubt.  The question should be treated as genuinely open.  And if the
field cannot bring itself to treat it as open---if the mere posing of
the question continues to be treated as a mark of incompetence rather
than as a legitimate scholarly inquiry---then the field has a problem
that no amount of consensus can resolve.

\vspace{2em}
\begin{center}
\rule{0.5\textwidth}{0.4pt}
\end{center}
\vspace{1em}

\noindent\textit{Acknowledgments.}  This paper benefited from extensive
critical dialogue conducted with Claude (Anthropic), which served as
adversarial interlocutor and peer reviewer throughout the analytical
process.  The arguments and conclusions are the author's own.

%% –– Bibliography ––
\bibliography{historicity}

\end{document}
